% Beamer slide deck for K. A. Novak,
% "Special Functions of Mathematical Physics: A Tourist's Guidebook"
% Chapter 7: Orthogonal Polynomials
% Self-contained slide text; local figures live in sfmp_slides/figs.

\documentclass[11pt]{beamer}

\usepackage[utf8]{inputenc}
\usepackage[T1]{fontenc}
\usepackage{lmodern}
\usepackage{amsmath,amssymb}
\usepackage{graphicx}

\usetheme{Madrid}

\author[]{\texorpdfstring{after K. A. Novak\\Compiled by Jake W. Liu}{after K. A. Novak, Compiled by Jake W. Liu}}
\title[Orthogonal Polynomials]{Chapter 7: Orthogonal Polynomials}
\subtitle{Special Functions of Mathematical Physics: A Tourist's Guidebook}

\setbeamertemplate{navigation symbols}{}
\setbeamertemplate{frametitle continuation}[from second]

\AtBeginSection[]
{
 \begin{frame}<beamer>
 \frametitle{Table of contents}
 \tableofcontents[currentsection]
 \end{frame}
}

\usepackage{Kyushu}

% Neutralize \index so source markup (if any) is harmless.
\makeatletter
\@ifundefined{index}{\newcommand{\index}[1]{}}{\renewcommand{\index}[1]{}}
\makeatother

\begin{document}

\newcounter{sfchapter}
\setcounter{sfchapter}{7}
\renewcommand{\thesection}{\arabic{sfchapter}.\arabic{section}}
\numberwithin{equation}{section}

\begin{frame}[c]\titlepage
\end{frame}

\begin{frame}{Outline}
\tableofcontents
\end{frame}

\begin{frame}{Where we are going}
\footnotesize
This chapter studies four families of \textbf{orthogonal polynomials}---Legendre,
Hermite, Laguerre, and Chebyshev---each the polynomial solution of a second-order
differential equation that arises by separating variables in a physical problem.
They appear in Laplace, Poisson, Helmholtz, and Schr\"odinger equations and their
time-dependent variants.
\begin{itemize}
  \item \textbf{Legendre} $P_n$: Laplace's equation $\nabla^2\psi=0$ in spherical
        coordinates (spherical harmonics).
  \item \textbf{Hermite} $H_n$: Schr\"odinger's equation for the quantum harmonic
        oscillator.
  \item \textbf{Laguerre}: generalized Laguerre polynomials arise in the radial
        hydrogenic equation; the standard family $L_n=L_n^{(0)}$ is developed here.
  \item \textbf{Chebyshev} $T_n$: defined through $\cos n\theta$; a natural interval
        basis, and tied to separated problems in elliptic coordinates.
\end{itemize}
Across the chapter we use the same toolkit: a defining ODE, a series solution when useful,
a \emph{generating function}, a \emph{recurrence}, and an \emph{orthogonality} relation.
For Legendre, Hermite, and Laguerre we also record Rodrigues' formulas and
Schl\"afli (contour) integrals.
\end{frame}

% ======================================================================
\section{Spherical Harmonics}
% ======================================================================

\begin{frame}[allowframebreaks]{Warm-up: 2D Laplace in polar coordinates}
\footnotesize

A function satisfying Laplace's equation $\nabla^2\psi=0$ is called a \emph{harmonic
function}. In two dimensions,
\begin{equation}
  \nabla^2\psi
   = \frac{1}{r}\,\frac{\partial}{\partial r}\!\left(r\,\frac{\partial\psi}{\partial r}\right)
   + \frac{1}{r^2}\,\frac{\partial^2\psi}{\partial\theta^2} = 0.
  \label{eq:laplace2d}
\end{equation}
Take a separated solution $\psi(r,\theta)=u(r)v(\theta)$. Then
\[
  \frac{\partial\psi}{\partial r}=u'(r)v(\theta),
  \qquad
  \frac{\partial^2\psi}{\partial\theta^2}=u(r)v''(\theta),
\]
so substitution into \eqref{eq:laplace2d} gives
\[
  \frac{v}{r}\frac{d}{dr}\!\left(ru'\right)+\frac{u}{r^2}v''=0.
\]
Multiplying by $r^2/(uv)$ yields
\[
  \frac{r(ru')'}{u}+\frac{v''}{v}=0,
\]
and therefore
\begin{equation}
  -\frac{v''}{v}=\frac{r(ru')'}{u}=\lambda,
  \label{eq:laplace2d-separation}
\end{equation}
because the left member depends only on $\theta$ and the right member only on $r$.

The angular equation is
\[
  v''+\lambda v=0.
\]
For $\lambda>0$, write $\lambda=\mu^2$; then
\[
  v(\theta)=A\cos(\mu\theta)+B\sin(\mu\theta).
\]
The periodicity condition $v(\theta+2\pi)=v(\theta)$ requires $\mu=n$ with
$n\in\mathbb Z$. Since the modes $n$ and $-n$ span the same real sine--cosine pair, we
may take $n\ge0$ and $\lambda=n^2$. If $\lambda=0$, then $v=A+B\theta$, and periodicity
forces $B=0$; this is the $n=0$ mode. If $\lambda<0$, the angular solutions are
exponentials rather than periodic functions, so no nonzero $2\pi$-periodic mode occurs.

The radial equation is
\[
  r(ru')'-n^2u=0.
\]
Expanding the derivative gives the Euler equation
\begin{equation}
  r^2u''+ru'-n^2u=0.
  \label{eq:laplace2d-radial}
\end{equation}
For $n\ge1$, try $u=r^p$. Since
\[
  u'=pr^{p-1},
  \qquad
  u''=p(p-1)r^{p-2},
\]
we obtain
\[
  r^2u''+ru'-n^2u
   =\bigl[p(p-1)+p-n^2\bigr]r^p
   =(p^2-n^2)r^p.
\]
Thus $p=\pm n$, and
\[
  u_n(r)=C_nr^n+D_nr^{-n}
  \qquad(n\ge1).
\]
For $n=0$, \eqref{eq:laplace2d-radial} becomes $(ru')'=0$. Integrating twice,
\[
  ru'=C_0,
  \qquad
  u'=\frac{C_0}{r},
  \qquad
  u(r)=D_0+C_0\log r.
\]

Consequently, on an annulus the general real separated expansion has the form
\begin{align*}
  \psi(r,\theta)
  &=A_0+B_0\log r \\
  &\quad+\sum_{n=1}^{\infty}
     \Big[(A_nr^n+B_nr^{-n})\cos(n\theta)
          +(C_nr^n+D_nr^{-n})\sin(n\theta)\Big].
\end{align*}
Inside a disk containing $r=0$, boundedness removes $\log r$ and all $r^{-n}$ terms:
\begin{equation}
  \psi(r,\theta)
   =A_0+\sum_{n=1}^{\infty}r^n
       \bigl[A_n\cos(n\theta)+C_n\sin(n\theta)\bigr].
  \label{eq:laplace2d-sol}
\end{equation}
Equivalently, with $z=re^{i\theta}$ and complex coefficients $c_n=A_n-iC_n$,
\[
  \psi(r,\theta)=\operatorname{Re}\!\left(\sum_{n=0}^{\infty}c_nz^n\right),
  \qquad z^n=r^ne^{in\theta}.
\]
Thus the regular real harmonic functions are real parts of analytic power series; exterior
or annular domains additionally allow negative powers and, for the angularly independent
$n=0$ mode, a logarithm.
\end{frame}

\begin{frame}[allowframebreaks]{3D Laplace in spherical coordinates}
\footnotesize

In spherical coordinates, Laplace's equation is
\begin{equation}
  \nabla^2\psi
  =
  \frac{1}{r^2}
    \frac{\partial}{\partial r}
      \left(r^2\frac{\partial\psi}{\partial r}\right)
  +\frac{1}{r^2\sin\theta}
    \frac{\partial}{\partial\theta}
      \left(\sin\theta\frac{\partial\psi}{\partial\theta}\right)
  +\frac{1}{r^2\sin^2\theta}
    \frac{\partial^2\psi}{\partial\phi^2}
  =0.
  \label{eq:laplace3d}
\end{equation}
First separate the radial variable:
\[
  \psi(r,\theta,\phi)=u(r)y(\theta,\phi).
\]
Then
\[
  \psi_r=u'y,
  \qquad
  \psi_\theta=uy_\theta,
  \qquad
  \psi_{\phi\phi}=uy_{\phi\phi}.
\]
Because $y$ is independent of $r$ and $u$ is independent of the angles,
\begin{align*}
  \frac{\partial}{\partial r}(r^2\psi_r)
    &=\frac{\partial}{\partial r}(r^2u'y)
      =y(r^2u')',\\
  \frac{\partial}{\partial\theta}(\sin\theta\,\psi_\theta)
    &=\frac{\partial}{\partial\theta}
      (\sin\theta\,u y_\theta)
      =u(\sin\theta\,y_\theta)_\theta.
\end{align*}
Substitution into \eqref{eq:laplace3d} gives
\[
  \frac{y}{r^2}(r^2u')'
  +\frac{u}{r^2\sin\theta}(\sin\theta\,y_\theta)_\theta
  +\frac{u}{r^2\sin^2\theta}y_{\phi\phi}=0.
\]
Multiplying by $r^2/(uy)$ gives
\begin{equation}
  \frac{(r^2u')'}{u}
  =
  -\frac{
      \dfrac{1}{\sin\theta}(\sin\theta\,y_\theta)_\theta
      +\dfrac{1}{\sin^2\theta}y_{\phi\phi}
    }{y}
  =\lambda.
  \label{eq:sep-radial3d}
\end{equation}
The first expression depends only on $r$ and the second only on
$(\theta,\phi)$, so each must equal a constant.

The radial equation is
\[
  (r^2u')'-\lambda u=0.
\]
Expanding the derivative,
\[
  (r^2u')'=r^2u''+2ru',
\]
so
\[
  r^2u''+2ru'-\lambda u=0.
\]
Try an Euler power $u=r^\sigma$.  Then
\[
  u'=\sigma r^{\sigma-1},
  \qquad
  u''=\sigma(\sigma-1)r^{\sigma-2},
\]
and therefore
\begin{align*}
  r^2u''+2ru'-\lambda u
   &=
    \left[
      \sigma(\sigma-1)+2\sigma-\lambda
    \right]r^\sigma\\
   &=
    \left[
      \sigma(\sigma+1)-\lambda
    \right]r^\sigma.
\end{align*}
The angular regularity calculation in the next frame gives
\[
  \lambda=n(n+1),
  \qquad n=0,1,2,\ldots.
\]
Hence
\begin{align*}
  \sigma(\sigma+1)-n(n+1)
   &=\sigma^2+\sigma-n^2-n\\
   &=(\sigma-n)(\sigma+n+1)=0.
\end{align*}
The two exponents are
\[
  \sigma=n,
  \qquad
  \sigma=-n-1,
\]
so
\begin{equation}
  u_n(r)=A_nr^n+B_nr^{-n-1}.
  \label{eq:spherical-radial-powers}
\end{equation}
The $r^n$ branch is regular at the origin, whereas the $r^{-n-1}$ branch
decays at infinity.

The remaining angular eigenproblem, before quantizing $\lambda$, is
\[
  \frac{1}{\sin\theta}
    \frac{\partial}{\partial\theta}
      \left(\sin\theta\,\frac{\partial y}{\partial\theta}\right)
  +\frac{1}{\sin^2\theta}\frac{\partial^2y}{\partial\phi^2}
  +\lambda y=0.
\]
The next frame separates $\theta$ and $\phi$, imposes single-valuedness, and
derives $\lambda=n(n+1)$.
\end{frame}

\begin{frame}[allowframebreaks]{The associated Legendre equation}
\footnotesize

Before imposing the angular eigenvalue, the angular equation obtained from
\eqref{eq:sep-radial3d} is
\[
  \frac{1}{\sin\theta}
    \frac{\partial}{\partial\theta}
      \left(\sin\theta\,\frac{\partial y}{\partial\theta}\right)
  +\frac{1}{\sin^2\theta}\frac{\partial^2y}{\partial\phi^2}
  +\lambda y=0.
\]
Set $y(\theta,\phi)=v(\theta)w(\phi)$.  Substitution gives
\[
  \frac{w}{\sin\theta}(\sin\theta\,v')'
  +\frac{v}{\sin^2\theta}w''
  +\lambda vw=0.
\]
Multiplication by $\sin^2\theta/(vw)$ yields
\[
  \frac{\sin\theta(\sin\theta\,v')'}{v}
  +\lambda\sin^2\theta
  +\frac{w''}{w}=0.
\]
The first two terms depend only on $\theta$, while the last depends only on
$\phi$.  Set
\[
  -\frac{w''}{w}=m^2.
\]
Then
\[
  w''+m^2w=0,
  \qquad
  w(\phi)=A\cos(m\phi)+B\sin(m\phi).
\]
Single-valuedness under $\phi\mapsto\phi+2\pi$ requires $m\in\mathbb Z$.
When $m=0$, the general solution $w=A+B\phi$ is periodic only if $B=0$.
Equivalently, all periodic modes may be written as $w=e^{im\phi}$.

The polar equation is therefore
\begin{equation}
  \sin\theta\,(\sin\theta\,v_\theta)_\theta
  +\lambda\sin^2\theta\,v-m^2v=0.
  \label{eq:theta-associated-legendre}
\end{equation}
Set $x=\cos\theta$ and regard $v$ as a function of $x$.  Since
\begin{equation}
  \frac{dx}{d\theta}=-\sin\theta,
  \qquad
  \frac{d}{d\theta}
    =-\sin\theta\frac{d}{dx},
  \label{eq:dthetadx}
\end{equation}
we have
\begin{align*}
  \frac{dv}{d\theta}
    &=\frac{dv}{dx}\frac{dx}{d\theta}
      =-\sin\theta\,v_x,\\
  \sin\theta\frac{dv}{d\theta}
    &=-\sin^2\theta\,v_x
      =-(1-x^2)v_x,\\
  \frac{d}{d\theta}
    \left(\sin\theta\frac{dv}{d\theta}\right)
    &=
      -\sin\theta\frac{d}{dx}
      \left[-(1-x^2)v_x\right]\\
    &=
      \sin\theta\frac{d}{dx}
      \left[(1-x^2)v_x\right].
\end{align*}
Multiplying the last expression by the outer $\sin\theta$ in
\eqref{eq:theta-associated-legendre} gives
\[
  \sin\theta(\sin\theta\,v_\theta)_\theta
    =(1-x^2)\frac{d}{dx}\left[(1-x^2)v_x\right].
\]
Hence
\[
  (1-x^2)\frac{d}{dx}\left[(1-x^2)v_x\right]
  +\lambda(1-x^2)v-m^2v=0.
\]
For $-1<x<1$, divide by $1-x^2$:
\begin{equation}
  \frac{d}{dx}\left[(1-x^2)v'(x)\right]
  +\left[
      \lambda-\frac{m^2}{1-x^2}
   \right]v(x)=0.
  \label{eq:assoc-legendre-general}
\end{equation}
After expanding the derivative,
\[
  (1-x^2)v''-2xv'
  +\left[
      \lambda-\frac{m^2}{1-x^2}
   \right]v=0.
\]

We now derive the allowed values of $\lambda$.  Put $p:=|m|$.  Near
$x=1$, let $z=1-x$, so
\[
  1-x^2=(1-x)(1+x)=z(2-z)\sim2z.
\]
If $v\sim z^\alpha$, then
\[
  v'=-\alpha z^{\alpha-1},
  \qquad
  v''=\alpha(\alpha-1)z^{\alpha-2}.
\]
The most singular terms in the expanded associated Legendre equation are
\begin{align*}
  (1-x^2)v''-2xv'-\frac{p^2}{1-x^2}v
   &\sim
    2z\,\alpha(\alpha-1)z^{\alpha-2}
    +2\alpha z^{\alpha-1}
    -\frac{p^2}{2z}z^\alpha\\
   &=
    \left(2\alpha^2-\frac{p^2}{2}\right)z^{\alpha-1}.
\end{align*}
The indicial equation is therefore
\[
  2\alpha^2-\frac{p^2}{2}=0,
  \qquad
  \alpha=\pm\frac p2.
\]
The regular branch uses $\alpha=p/2$; the same endpoint behavior occurs at
$x=-1$.  This motivates
\[
  s:=1-x^2,
  \qquad
  v(x)=s^{p/2}q(x).
\]
Let $f(x)=s^{p/2}$.  Since $s'=-2x$,
\begin{align*}
  f'
    &=\frac p2s^{p/2-1}s'
      =-pxs^{p/2-1},\\
  f''
    &=-p s^{p/2-1}
      +p(p-2)x^2s^{p/2-2}.
\end{align*}
Also,
\[
  v'=f'q+fq',
  \qquad
  v''=f''q+2f'q'+fq''.
\]
Substitute these expressions into the expanded form of
\eqref{eq:assoc-legendre-general}:
\begin{align*}
  0
   &=
    s(f''q+2f'q'+fq'')
    -2x(f'q+fq')
    +\left(\lambda-\frac{p^2}{s}\right)fq\\
   &=
    f\Bigg\{
      s q''
      +\left(2s\frac{f'}{f}-2x\right)q'
      +\left[
        s\frac{f''}{f}
        -2x\frac{f'}{f}
        +\lambda-\frac{p^2}{s}
       \right]q
    \Bigg\}.
\end{align*}
The coefficient of $q'$ is
\[
  2s\frac{f'}{f}-2x
    =2s\left(-\frac{px}{s}\right)-2x
    =-2(p+1)x.
\]
For the coefficient of $q$,
\begin{align*}
  s\frac{f''}{f}
  -2x\frac{f'}{f}
  +\lambda-\frac{p^2}{s}
   &=
    \left[-p+\frac{p(p-2)x^2}{s}\right]
    +\frac{2px^2}{s}
    +\lambda-\frac{p^2}{s}\\
   &=
    \lambda-p
    +\frac{p^2x^2-p^2}{s}\\
   &=
    \lambda-p
    +p^2\frac{x^2-1}{1-x^2}\\
   &=
    \lambda-p-p^2\\
   &=
    \lambda-p(p+1).
\end{align*}
Thus $q$ satisfies
\begin{equation}
  (1-x^2)q''-2(p+1)xq'
  +\bigl[\lambda-p(p+1)\bigr]q=0.
  \label{eq:assoc-reduced}
\end{equation}

Let
\[
  q(x)=\sum_{j=0}^{\infty}a_jx^j.
\]
Substitution into \eqref{eq:assoc-reduced}, followed by the same coefficient
shifts used for the ordinary Legendre equation, gives
\[
  0
   =
  \sum_{j=0}^{\infty}
  \left[
    (j+2)(j+1)a_{j+2}+C_j a_j
  \right]x^j,
\]
where
\begin{align*}
  C_j
   &=
    \lambda-j(j-1)-2(p+1)j-p(p+1)\\
   &=
    \lambda-\left[j^2+(2p+1)j+p(p+1)\right]\\
   &=
    \lambda-(j+p)(j+p+1).
\end{align*}
Therefore
\begin{equation}
  a_{j+2}
   =
  -\frac{\lambda-(j+p)(j+p+1)}
         {(j+2)(j+1)}a_j.
  \label{eq:assoc-coeff-recurrence}
\end{equation}
The even and odd coefficient chains are independent.  A polynomial of degree
$\ell\ge0$ is obtained by retaining the parity chain containing $x^\ell$,
setting the other initial coefficient to zero, and requiring the numerator to
vanish at $j=\ell$:
\[
  \lambda=(\ell+p)(\ell+p+1).
\]
Define
\[
  n:=\ell+p.
\]
Then
\[
  \lambda=n(n+1),
  \qquad
  \ell=n-p,
  \qquad
  0\le p=|m|\le n.
\]
This derives both the angular eigenvalue and the index restriction.

With $\lambda=n(n+1)$, \eqref{eq:assoc-legendre-general} becomes
\begin{equation}
  \frac{d}{dx}\left[(1-x^2)v'(x)\right]
  +\left[
      n(n+1)-\frac{m^2}{1-x^2}
   \right]v(x)=0.
  \label{eq:assoc-legendre}
\end{equation}

To connect the solution with the ordinary Legendre polynomial, differentiate
\eqref{eq:legendre-eq-expanded} $p$ times.  Since derivatives of
$1-x^2$ above order $2$ vanish, Leibniz' rule gives
\begin{align*}
  \frac{d^p}{dx^p}\left[(1-x^2)P_n''\right]
   &=
    (1-x^2)P_n^{(p+2)}
    +p(-2x)P_n^{(p+1)}
    +\binom{p}{2}(-2)P_n^{(p)}\\
   &=
    (1-x^2)P_n^{(p+2)}
    -2pxP_n^{(p+1)}
    -p(p-1)P_n^{(p)},\\
  \frac{d^p}{dx^p}\left[-2xP_n'\right]
   &=
    -2xP_n^{(p+1)}-2pP_n^{(p)}.
\end{align*}
Therefore $q=P_n^{(p)}$ satisfies
\[
  (1-x^2)q''-2(p+1)xq'
  +\left[n(n+1)-p(p+1)\right]q=0,
\]
which is exactly \eqref{eq:assoc-reduced} with $\lambda=n(n+1)$.  Multiplying
by the regular endpoint factor $(1-x^2)^{p/2}$ therefore gives an associated
Legendre solution.  A standard phase convention defines
\begin{equation}
  P_n^p(x)
    :=(-1)^p(1-x^2)^{p/2}
      \frac{d^pP_n(x)}{dx^p},
  \qquad p=0,1,\ldots,n.
  \label{eq:associated-legendre-definition}
\end{equation}
Because the differential equation depends on $m^2$, an unnormalized angular
basis for both signs of $m$ is
\[
  \mathcal Y_n^m(\theta,\phi)
    =P_n^{|m|}(\cos\theta)e^{im\phi},
  \qquad -n\le m\le n.
\]
Normalized spherical harmonics multiply this expression by a normalization
constant and adopt a fixed phase convention.

Combining radial and angular factors gives the regular interior expansion
\begin{equation}
  \psi_{\mathrm{int}}(r,\theta,\phi)
    =\sum_{n=0}^{\infty}\sum_{m=-n}^{n}
      a_{nm}r^n\mathcal Y_n^m(\theta,\phi),
  \label{eq:laplace3d-sol}
\end{equation}
and the decaying exterior expansion
\[
  \psi_{\mathrm{ext}}(r,\theta,\phi)
    =\sum_{n=0}^{\infty}\sum_{m=-n}^{n}
      b_{nm}r^{-n-1}\mathcal Y_n^m(\theta,\phi).
\]
On a spherical shell, both radial branches may be present.
\end{frame}

\begin{frame}{Low spherical-harmonic nodal patterns}
\footnotesize
\begin{center}
  \includegraphics[width=0.88\linewidth]{figs/ch7_spherical_harmonics.pdf}
\end{center}
\scriptsize
Rows show degree $n$ and columns show order $m$.  Let $p=|m|$.  Sturm
oscillation for the associated Legendre equation gives $n-p$ zeros of
$P_n^p(\cos\theta)$ for $0<\theta<\pi$, hence $n-p$ nodal latitude circles.
For a real mode, the azimuthal factor is $\cos(p\phi)$ or $\sin(p\phi)$.
It has $2p$ zero rays in $0\le\phi<2\pi$; opposite rays form one plane through
the polar axis, so there are $p$ nodal meridian planes.
\end{frame}

\begin{frame}{Real-part spherical-harmonic surfaces}
\footnotesize
\begin{center}
  \includegraphics[width=0.96\linewidth]{figs/ch7_spherical_harmonics_surfaces.pdf}
\end{center}
The book's surface picture emphasizes amplitude and sign, not only nodal counts. Blue and
red indicate opposite signs of \(\operatorname{Re}Y_n^m\); the lobes shrink to zero on the
nodal latitude and meridian curves.
\end{frame}

\begin{frame}[allowframebreaks]{Axisymmetry: the ordinary Legendre equation}
\footnotesize

If the problem is axisymmetric about the polar axis, then nothing depends on $\phi$, so
$m=0$.  This is weaker than full radial symmetry: full radial symmetry leaves only the
$n=0$ mode, whereas axisymmetry permits every degree $n$ with $m=0$.

Setting $m=0$ in \eqref{eq:assoc-legendre} gives
\begin{equation}
  \frac{d}{dx}\!\left[(1-x^2)v'(x)\right]+n(n+1)v(x)=0.
  \label{eq:legendre-eq}
\end{equation}
Using the product rule,
\[
  \frac{d}{dx}\!\left[(1-x^2)v'\right]
   =(1-x^2)v''+\frac{d}{dx}(1-x^2)v'
   =(1-x^2)v''-2xv',
\]
so the equivalent expanded equation is
\begin{equation}
  (1-x^2)v''-2xv'+n(n+1)v=0.
  \label{eq:legendre-eq-expanded}
\end{equation}

To see why integer $n$ produces a polynomial, set
\[
  v(x)=\sum_{j=0}^{\infty}a_jx^j.
\]
Then
\begin{align*}
  v'
   &=\sum_{j=0}^{\infty}(j+1)a_{j+1}x^j,\\
  v''
   &=\sum_{j=0}^{\infty}(j+2)(j+1)a_{j+2}x^j,\\
  -x^2v''
   &=-\sum_{j=2}^{\infty}j(j-1)a_jx^j
    =-\sum_{j=0}^{\infty}j(j-1)a_jx^j,\\
  -2xv'
   &=-2\sum_{j=1}^{\infty}j\,a_jx^j
    =-2\sum_{j=0}^{\infty}j\,a_jx^j.
\end{align*}
Substitution into \eqref{eq:legendre-eq-expanded} gives
\begin{align*}
  0
   &=
    \sum_{j=0}^{\infty}
    \left[
      (j+2)(j+1)a_{j+2}
      -j(j-1)a_j-2ja_j+n(n+1)a_j
    \right]x^j\\
   &=
    \sum_{j=0}^{\infty}
    \left[
      (j+2)(j+1)a_{j+2}
      +\bigl(n(n+1)-j(j+1)\bigr)a_j
    \right]x^j.
\end{align*}
Since
\[
  n(n+1)-j(j+1)
   =n^2+n-j^2-j
   =(n-j)(n+j+1),
\]
the coefficient recurrence is
\begin{equation}
  a_{j+2}
   =-\frac{(n-j)(n+j+1)}{(j+2)(j+1)}\,a_j.
  \label{eq:legendre-coeff-recurrence}
\end{equation}
This recurrence separates the even and odd coefficients.  For integer $n\ge0$, the
chain having the same parity as $n$ reaches $j=n$, where
\[
  a_{n+2}
   =-\frac{(n-n)(n+n+1)}{(n+2)(n+1)}a_n
   =0.
\]
Every later coefficient in that chain is consequently zero.  Setting the
opposite-parity initial coefficient to zero leaves a polynomial of degree $n$.
The remaining multiplicative constant is fixed by the standard normalization
\[
  P_n(1)=1.
\]
Thus $P_n$ contains only the powers $x^n,x^{n-2},\ldots$ and obeys
\[
  P_n(-x)=(-1)^nP_n(x).
\]
The generating function below determines all normalized coefficients explicitly.
\end{frame}

% ======================================================================
\section{Legendre Polynomial}
% ======================================================================

\begin{frame}[allowframebreaks]{Generating function from the Coulomb potential}
\footnotesize

Place a point charge one unit from the origin on the positive polar axis. For a field point
with spherical coordinates $(r,\theta,\phi)$, the law of cosines gives
\[
  R^2=r^2+1^2-2r(1)\cos\theta
      =1-2r\cos\theta+r^2,
\]
so the Coulomb potential is
\begin{equation}
  \varphi(r,\theta)=\frac{C}{R}
  =\frac{C}{\sqrt{1-2r\cos\theta+r^2}}.
  \label{eq:coulomb}
\end{equation}
\begin{center}
  \includegraphics[width=0.42\linewidth]{figs/legendre_coulomb_geometry.pdf}
\end{center}
For $0\le r<1$, the charge lies outside the ball containing the field point; hence
$\varphi$ is harmonic there. The source and geometry are axisymmetric, so the regular
interior expansion has the form
\[
  \varphi(r,\theta)
  =\sum_{n=0}^{\infty}c_n r^nP_n(\cos\theta).
\]
Dividing by the nonzero scale factor $C$ gives
\begin{equation}
  \frac{1}{\sqrt{1-2r\cos\theta+r^2}}
  =\sum_{n=0}^{\infty}a_n r^nP_n(\cos\theta),
  \qquad a_n=\frac{c_n}{C}.
  \label{eq:coulomb-expand}
\end{equation}

Set $\theta=0$. Then $\cos\theta=1$, $P_n(1)=1$, and
\[
  1-2r\cos0+r^2=1-2r+r^2=(1-r)^2.
\]
Because $0\le r<1$, $|1-r|=1-r$, so \eqref{eq:coulomb-expand} becomes
\[
  \frac{1}{1-r}=\sum_{n=0}^{\infty}a_nr^n.
\]
The geometric series is
\[
  \frac{1}{1-r}=1+r+r^2+r^3+\cdots
  =\sum_{n=0}^{\infty}r^n,
  \qquad |r|<1.
\]
Uniqueness of power-series coefficients therefore gives
\[
  a_n=1\qquad(n=0,1,2,\ldots).
\]
Substitution into \eqref{eq:coulomb-expand} yields
\[
  \frac{1}{\sqrt{1-2r\cos\theta+r^2}}
  =\sum_{n=0}^{\infty}P_n(\cos\theta)r^n.
\]
Writing $t=r$ and $x=\cos\theta$ gives the Legendre generating function
\begin{equation}
  g(x,t)=\frac{1}{\sqrt{1-2xt+t^2}}
  =\sum_{n=0}^{\infty}P_n(x)t^n,
  \qquad |t|<1,\quad -1\le x\le1.
  \label{eq:legendre-gen}
\end{equation}
At $t=0$, this formula gives $P_0(x)=1$. Differentiating once and setting $t=0$ gives
$P_1(x)=x$.
\end{frame}

\begin{frame}[allowframebreaks]{Three-term recurrence}
\footnotesize

Let
\[
  D(x,t)=1-2xt+t^2,
  \qquad
  g(x,t)=D(x,t)^{-1/2}.
\]
Differentiating with respect to $t$ gives
\begin{align*}
  \frac{\partial g}{\partial t}
  &=-\frac12D^{-3/2}\frac{\partial D}{\partial t} \\
  &=-\frac12D^{-3/2}(-2x+2t) \\
  &=\frac{x-t}{D^{3/2}}.
\end{align*}
Termwise differentiation of the power series gives
\begin{equation}
  \frac{\partial g}{\partial t}
  =\sum_{n=1}^{\infty}nP_n(x)t^{n-1}.
  \label{eq:dgdt}
\end{equation}
Multiplying $g_t=(x-t)D^{-3/2}$ by $D$ gives
\[
  (x-t)g=Dg_t.
\]
Substituting both power series,
\[
  (x-t)\sum_{n=0}^{\infty}P_nt^n
  =(1-2xt+t^2)\sum_{n=1}^{\infty}nP_nt^{n-1},
\]
where $P_n=P_n(x)$. Expanding both sides,
\begin{align}
  \sum_{n=0}^{\infty}xP_nt^n
  -\sum_{n=0}^{\infty}P_nt^{n+1}
  &=\sum_{n=1}^{\infty}nP_nt^{n-1}
    -2\sum_{n=1}^{\infty}nxP_nt^n
    +\sum_{n=1}^{\infty}nP_nt^{n+1}.
  \label{eq:distribute}
\end{align}

We now compute the coefficient of a fixed power $t^n$.
On the left,
\[
  [t^n]\left(\sum_{j=0}^{\infty}xP_jt^j\right)=xP_n,
\]
and
\[
  [t^n]\left(\sum_{j=0}^{\infty}P_jt^{j+1}\right)
  =\begin{cases}
      0,&n=0,\\
      P_{n-1},&n\ge1.
    \end{cases}
\]
It is convenient to set $P_{-1}=0$, so the left coefficient is $xP_n-P_{n-1}$ for all
$n\ge0$.

On the right,
\[
  [t^n]\left(\sum_{j=1}^{\infty}jP_jt^{j-1}\right)
  =(n+1)P_{n+1},
\]
\[
  [t^n]\left(-2\sum_{j=1}^{\infty}jxP_jt^j\right)
  =-2nxP_n,
\]
and
\[
  [t^n]\left(\sum_{j=1}^{\infty}jP_jt^{j+1}\right)
  =\begin{cases}
     0,&n=0,1,\\
     (n-1)P_{n-1},&n\ge2.
   \end{cases}
\]
The last expression is also $(n-1)P_{n-1}$ for all $n\ge0$ under $P_{-1}=0$.
Equating coefficients therefore gives
\[
  xP_n-P_{n-1}
  =(n+1)P_{n+1}-2nxP_n+(n-1)P_{n-1}.
\]
Move the terms containing $P_n$ and $P_{n-1}$ to the same side:
\[
  (2n+1)xP_n-nP_{n-1}=(n+1)P_{n+1}.
\]
Thus
\begin{equation}
  (n+1)P_{n+1}(x)
  =(2n+1)xP_n(x)-nP_{n-1}(x),
  \qquad n\ge0,
  \label{eq:legendre-recur}
\end{equation}
with $P_{-1}=0$ and $P_0=1$. For $n=0$, the recurrence gives
$P_1=xP_0=x$, consistent with the generating function.
\end{frame}

\begin{frame}[allowframebreaks]{First few Legendre polynomials}
\footnotesize

The initial values are
\[
  P_0(x)=1,
  \qquad
  P_1(x)=x.
\]
For $n=1$, the recurrence \eqref{eq:legendre-recur} gives
\begin{align*}
  (1+1)P_2
   &=(2\cdot1+1)xP_1-1\cdot P_0,\\
  2P_2
   &=3x(x)-1,\\
  2P_2
   &=3x^2-1,
\end{align*}
so
\[
  P_2(x)=\frac12(3x^2-1).
\]
For $n=2$,
\begin{align*}
  (2+1)P_3
   &=(2\cdot2+1)xP_2-2P_1,\\
  3P_3
   &=5x\left[\frac12(3x^2-1)\right]-2x,\\
  3P_3
   &=\frac12(15x^3-5x)-\frac{4x}{2},\\
  3P_3
   &=\frac12(15x^3-9x),\\
  P_3
   &=\frac16(15x^3-9x)
    =\frac12(5x^3-3x).
\end{align*}
Thus
\begin{align*}
  P_0(x)&=1,&
  P_1(x)&=x,\\[2pt]
  P_2(x)&=\frac12(3x^2-1),&
  P_3(x)&=\frac12(5x^3-3x).
\end{align*}
As checks, $P_n(1)=1$ in all four cases, and $P_n(-x)=(-1)^nP_n(x)$.
\end{frame}

\begin{frame}{Shapes of the first Legendre polynomials}
\footnotesize
\begin{center}
  \includegraphics[width=0.86\linewidth]{figs/legendre-P.pdf}
\end{center}
Each $P_n$ is bounded on $[-1,1]$ and normalized by $P_n(1)=1$. Increasing $n$ adds
interior zeros while preserving the Sturm--Liouville orthogonality on the same interval.
\end{frame}

\begin{frame}[allowframebreaks]{Contour integral from the generating function}
\footnotesize

For fixed $x\in[-1,1]$, choose the branch of
\[
  g(x,t)=(1-2xt+t^2)^{-1/2}
\]
that satisfies $g(x,0)=1$. It is analytic in a neighborhood of $t=0$, and
\[
  g(x,t)=\sum_{n=0}^{\infty}P_n(x)t^n.
\]
Cauchy's coefficient formula states that, for any positively oriented simple contour
$\gamma$ around $t=0$ lying inside the domain of analyticity,
\[
  [t^m]g(x,t)=\frac{1}{2\pi i}\oint_\gamma\frac{g(x,t)}{t^{m+1}}\,dt.
\]
Here $[t^m]g=P_m(x)$, so
\begin{align*}
  \oint_\gamma\frac{dt}{t^{m+1}\sqrt{1-2xt+t^2}}
   &=\oint_\gamma\sum_{n=0}^{\infty}P_n(x)t^{n-m-1}\,dt\\
   &=\sum_{n=0}^{\infty}P_n(x)
       \oint_\gamma t^{n-m-1}\,dt\\
   &=P_m(x)(2\pi i),
\end{align*}
because
\[
  \oint_\gamma t^k\,dt
  =\begin{cases}
      2\pi i,&k=-1,\\
      0,&k\ne-1.
    \end{cases}
\]
Therefore
\begin{equation}
  P_m(x)=\frac{1}{2\pi i}
  \oint_\gamma\frac{dt}{t^{m+1}\sqrt{1-2xt+t^2}}.
  \label{eq:legendre-contour}
\end{equation}

For $-1<x<1$, the zeros of the radicand are
\begin{align*}
  t_{\pm}
   &=\frac{2x\pm\sqrt{4x^2-4}}{2}\\
   &=x\pm i\sqrt{1-x^2}.
\end{align*}
Writing $x=\cos\theta$ gives
\[
  t_\pm=\cos\theta\pm i\sin\theta=e^{\pm i\theta},
\]
so both branch points lie on $|t|=1$. A sufficiently small circle $|t|=\rho<1$ is
therefore an admissible contour. At $x=\pm1$ the two roots coalesce and the formula is
understood by the corresponding analytic limiting branch.
\end{frame}

\begin{frame}[allowframebreaks]{Explicit Legendre series from the generating function}
\footnotesize

Starting from \eqref{eq:legendre-gen}, use the generalized binomial expansion
\[
  (1+u)^{-1/2}
    =\sum_{q=0}^{\infty}\binom{-1/2}{q}u^q.
\]
For $q\ge0$,
\begin{align*}
  \binom{-1/2}{q}
   &=
    \frac{(-1/2)(-3/2)\cdots(-(2q-1)/2)}{q!}\\
   &=
    (-1)^q\frac{1\cdot3\cdot5\cdots(2q-1)}{2^q q!}\\
   &=
    (-1)^q\frac{(2q)!}{4^q(q!)^2},
\end{align*}
where the last equality uses
\[
  (2q)!
    =(2\cdot4\cdots2q)(1\cdot3\cdots(2q-1))
    =2^q q!(1\cdot3\cdots(2q-1)).
\]
With $u=-2xt+t^2$,
\begin{align*}
  g(x,t)
   &=(1-2xt+t^2)^{-1/2}\\
   &=\sum_{q=0}^{\infty}
      \binom{-1/2}{q}(-2xt+t^2)^q.
\end{align*}
Expand the $q$th power:
\begin{align*}
  (-2xt+t^2)^q
   &=
    \sum_{k=0}^{q}
      \binom{q}{k}(-2xt)^{q-k}(t^2)^k\\
   &=
    \sum_{k=0}^{q}
      \binom{q}{k}(-2x)^{q-k}t^{q+k}.
\end{align*}
A term contributes to the coefficient of $t^n$ precisely when
\[
  q+k=n,
  \qquad\text{so}\qquad
  q=n-k.
\]
The condition $0\le k\le q$ becomes
\[
  0\le k\le n-k,
  \qquad\text{or}\qquad
  0\le k\le\left\lfloor\frac n2\right\rfloor.
\]
Consequently,
\begin{align*}
  P_n(x)
   &=[t^n]\,g(x,t)\\
   &=
    \sum_{k=0}^{\lfloor n/2\rfloor}
      \binom{-1/2}{n-k}
      \binom{n-k}{k}
      (-2x)^{n-2k}\\
   &=
    \sum_{k=0}^{\lfloor n/2\rfloor}
      \left[
        (-1)^{n-k}
        \frac{(2n-2k)!}
             {4^{n-k}((n-k)!)^2}
      \right]
      \left[
        \frac{(n-k)!}{k!(n-2k)!}
      \right]
      (-1)^{n-2k}2^{n-2k}x^{n-2k}.
\end{align*}
The signs and powers of $2$ simplify as
\[
  (-1)^{n-k+n-2k}=(-1)^{2n-3k}=(-1)^k,
  \qquad
  \frac{2^{n-2k}}{4^{n-k}}
    =\frac{2^{n-2k}}{2^{2n-2k}}
    =2^{-n}.
\]
Therefore
\begin{equation}
  P_n(x)
   =\sum_{k=0}^{\lfloor n/2\rfloor}
      (-1)^k
      \frac{(2n-2k)!}
           {2^n k!(n-k)!(n-2k)!}
      x^{n-2k}.
  \label{eq:legendre-series}
\end{equation}
This coefficient formula has the parity predicted by
\eqref{eq:legendre-coeff-recurrence}.  Its normalization follows directly from
the generating function: at $x=1$,
\[
  \sum_{n=0}^{\infty}P_n(1)t^n
   =\frac{1}{\sqrt{1-2t+t^2}}
   =\frac{1}{1-t}
   =\sum_{n=0}^{\infty}t^n
   \qquad (|t|<1),
\]
so $P_n(1)=1$ for every $n$.
\end{frame}

\begin{frame}[allowframebreaks]{Rodrigues' formula and Schl\"afli integral}
\footnotesize

Expand $(x^2-1)^n$ by the binomial theorem:
\begin{equation}
  (x^2-1)^n
   =\sum_{j=0}^{n}
      \binom{n}{j}(x^2)^j(-1)^{n-j}
   =\sum_{j=0}^{n}
      (-1)^{n-j}\frac{n!}{j!(n-j)!}x^{2j}.
  \label{eq:binom}
\end{equation}
For each monomial,
\[
  \frac{d^n}{dx^n}x^{2j}
  =
  \begin{cases}
    \dfrac{(2j)!}{(2j-n)!}x^{2j-n}, & 2j\ge n,\\[0.7em]
    0, & 2j<n.
  \end{cases}
\]
Thus only $j\ge\lceil n/2\rceil$ survives:
\[
  \frac{d^n}{dx^n}(x^2-1)^n
   =
    \sum_{j=\lceil n/2\rceil}^{n}
      (-1)^{n-j}
      \frac{n!}{j!(n-j)!}
      \frac{(2j)!}{(2j-n)!}x^{2j-n}.
\]
Set $k=n-j$.  Then
\[
  j=n-k,
  \qquad
  0\le k\le\left\lfloor\frac n2\right\rfloor,
  \qquad
  2j-n=n-2k,
  \qquad
  (-1)^{n-j}=(-1)^k.
\]
The factorials transform as
\[
  j!=(n-k)!,
  \quad
  (n-j)!=k!,
  \quad
  (2j)!=(2n-2k)!,
  \quad
  (2j-n)!=(n-2k)!.
\]
Therefore
\begin{align*}
  \frac{d^n}{dx^n}(x^2-1)^n
   &=
    n!\sum_{k=0}^{\lfloor n/2\rfloor}
      (-1)^k
      \frac{(2n-2k)!}
           {k!(n-k)!(n-2k)!}
      x^{n-2k}.
\end{align*}
Comparison with \eqref{eq:legendre-series} shows that
\[
  \frac{d^n}{dx^n}(x^2-1)^n=2^n n!P_n(x).
\]
Hence Rodrigues' formula is
\begin{equation}
  P_n(x)
   =\frac{1}{2^n n!}
      \frac{d^n}{dx^n}(x^2-1)^n.
  \label{eq:rodrigues-legendre}
\end{equation}

Cauchy's differentiation formula states that
\[
  f^{(n)}(x)
    =\frac{n!}{2\pi i}
      \oint_\Gamma\frac{f(t)}{(t-x)^{n+1}}\,dt,
\]
where $\Gamma$ is a positively oriented contour enclosing $t=x$.  Apply it to
$f(t)=(t^2-1)^n$:
\begin{align*}
  P_n(x)
   &=
    \frac{1}{2^n n!}f^{(n)}(x)\\
   &=
    \frac{1}{2^n n!}
    \frac{n!}{2\pi i}
    \oint_\Gamma
      \frac{(t^2-1)^n}{(t-x)^{n+1}}\,dt\\
   &=
    \frac{1}{2^{n+1}\pi i}
    \oint_\Gamma
      \frac{(t^2-1)^n}{(t-x)^{n+1}}\,dt.
\end{align*}
Thus the Schl\"afli representation is
\begin{equation}
  P_n(x)
    =\frac{1}{2^{n+1}\pi i}
      \oint_\Gamma
      \frac{(t^2-1)^n}{(t-x)^{n+1}}\,dt.
  \label{eq:schlafli-legendre}
\end{equation}
For integer $n$, the numerator is a polynomial, so the only singularity of the
integrand enclosed by a sufficiently small $\Gamma$ is the pole at $t=x$.
\end{frame}

\begin{frame}[allowframebreaks]{Orthogonality and norm of Legendre polynomials}
\footnotesize

The Legendre equation is
\begin{equation}
  \frac{d}{dx}\!\left[(1-x^2)P_n'(x)\right]
  +n(n+1)P_n(x)=0.
  \label{eq:legendre-sl}
\end{equation}
Write the equations for $P_n$ and $P_m$:
\begin{align*}
  \bigl[(1-x^2)P_n'\bigr]'+n(n+1)P_n&=0,\\
  \bigl[(1-x^2)P_m'\bigr]'+m(m+1)P_m&=0.
\end{align*}
Multiply the first by $P_m$, the second by $P_n$, and subtract:
\[
  P_m\bigl[(1-x^2)P_n'\bigr]'
  -P_n\bigl[(1-x^2)P_m'\bigr]'
  +\bigl[n(n+1)-m(m+1)\bigr]P_nP_m=0.
\]
The first two terms are a total derivative because
\[
  \frac{d}{dx}\!\left[(1-x^2)(P_mP_n'-P_nP_m')\right]
  =P_m\bigl[(1-x^2)P_n'\bigr]'
   -P_n\bigl[(1-x^2)P_m'\bigr]'.
\]
Integrating from $-1$ to $1$ gives
\begin{align*}
  &\left[(1-x^2)(P_mP_n'-P_nP_m')\right]_{-1}^{1}\\
  &\qquad+\bigl[n(n+1)-m(m+1)\bigr]
    \int_{-1}^{1}P_n(x)P_m(x)\,dx=0.
\end{align*}
The boundary term vanishes because $1-x^2=0$ at $x=\pm1$ and the polynomials and their
derivatives are finite there. If $n\ne m$, then
\[
  n(n+1)-m(m+1)=(n-m)(n+m+1)\ne0,
\]
so
\begin{equation}
  \int_{-1}^{1}P_n(x)P_m(x)\,dx=0
  \qquad(n\ne m).
  \label{eq:legendre-orthogonal-zero}
\end{equation}

It remains to compute the norm. Rodrigues' formula gives
\[
  I_n:=\int_{-1}^{1}P_n(x)^2\,dx
  =\frac{1}{2^n n!}\int_{-1}^{1}
    P_n(x)\frac{d^n}{dx^n}(x^2-1)^n\,dx.
\]
Integrate by parts once:
\begin{align*}
  I_n
  &=\frac{1}{2^n n!}
    \left[P_n\frac{d^{n-1}}{dx^{n-1}}(x^2-1)^n\right]_{-1}^{1}\\
  &\quad-\frac{1}{2^n n!}\int_{-1}^{1}
    P_n'\frac{d^{n-1}}{dx^{n-1}}(x^2-1)^n\,dx.
\end{align*}
Because $(x^2-1)^n=(x-1)^n(x+1)^n$ has a zero of order $n$ at each endpoint,
\[
  \left.\frac{d^j}{dx^j}(x^2-1)^n\right|_{x=\pm1}=0,
  \qquad 0\le j\le n-1.
\]
Thus every boundary term in $n$ successive integrations by parts vanishes, and
\begin{equation}
  I_n=\frac{(-1)^n}{2^n n!}\int_{-1}^{1}
  P_n^{(n)}(x)(x^2-1)^n\,dx.
  \label{eq:legendre-norm-ibp}
\end{equation}

The leading term of $P_n$ comes from $k=0$ in \eqref{eq:legendre-series}:
\[
  P_n(x)=\frac{(2n)!}{2^n(n!)^2}x^n+\text{terms of degree at most }n-2.
\]
After $n$ derivatives, all lower-degree terms vanish and
\[
  P_n^{(n)}(x)
  =n!\frac{(2n)!}{2^n(n!)^2}
  =\frac{(2n)!}{2^n n!}.
\]
Also $(x^2-1)^n=(-1)^n(1-x^2)^n$. Substitution into
\eqref{eq:legendre-norm-ibp} gives
\[
  I_n=\frac{(2n)!}{2^{2n}(n!)^2}
  \int_{-1}^{1}(1-x^2)^n\,dx.
\]
To evaluate the remaining integral, set $x=2u-1$. Then
\[
  dx=2\,du,
  \qquad
  1-x^2=1-(2u-1)^2=4u(1-u),
\]
with $x=-1\leftrightarrow u=0$ and $x=1\leftrightarrow u=1$. Hence
\begin{align*}
  \int_{-1}^{1}(1-x^2)^n\,dx
  &=2\cdot4^n\int_0^1u^n(1-u)^n\,du\\
  &=2^{2n+1}B(n+1,n+1)\\
  &=2^{2n+1}\frac{\Gamma(n+1)^2}{\Gamma(2n+2)}\\
  &=2^{2n+1}\frac{(n!)^2}{(2n+1)!}.
\end{align*}
Therefore
\begin{align*}
  I_n
  &=\frac{(2n)!}{2^{2n}(n!)^2}
    \frac{2^{2n+1}(n!)^2}{(2n+1)!}\\
  &=2\frac{(2n)!}{(2n+1)(2n)!}\\
  &=\frac{2}{2n+1}.
\end{align*}
Combining the orthogonality and norm,
\begin{equation}
  \int_{-1}^{1}P_n(x)P_m(x)\,dx
  =\frac{2}{2n+1}\,\delta_{nm}.
  \label{eq:legendre-orthogonality-full}
\end{equation}
\end{frame}

\begin{frame}[allowframebreaks]{Fourier--Legendre series}
\footnotesize

The Legendre polynomials are orthogonal and complete in $L^2[-1,1]$. Thus, for
$f\in L^2[-1,1]$,
\[
  f(x)=\sum_{n=0}^{\infty}a_nP_n(x)
\]
in the $L^2$ sense. Multiply by $P_m(x)$ and integrate:
\begin{align*}
  \int_{-1}^{1}f(x)P_m(x)\,dx
  &=\int_{-1}^{1}
    \left(\sum_{n=0}^{\infty}a_nP_n(x)\right)P_m(x)\,dx\\
  &=\sum_{n=0}^{\infty}a_n
    \int_{-1}^{1}P_n(x)P_m(x)\,dx\\
  &=\sum_{n=0}^{\infty}a_n
    \frac{2}{2n+1}\delta_{nm}\\
  &=a_m\frac{2}{2m+1}.
\end{align*}
Solving for $a_m$ gives
\begin{equation}
  a_m=\frac{2m+1}{2}\int_{-1}^{1}f(x)P_m(x)\,dx.
  \label{eq:fourier-legendre-coeff}
\end{equation}
Therefore the Fourier--Legendre expansion is
\begin{equation}
  f(x)=\sum_{n=0}^{\infty}
  \left[\frac{2n+1}{2}\int_{-1}^{1}f(s)P_n(s)\,ds\right]P_n(x).
  \label{eq:fourier-legendre}
\end{equation}
The dummy integration variable $s$ distinguishes the coefficient integral from the point
$x$ at which the series is evaluated. Additional regularity assumptions on $f$ provide
stronger pointwise or uniform convergence statements; completeness alone guarantees the
$L^2$ expansion.
\end{frame}

\begin{frame}{Useful formulas for the Legendre polynomial}
\footnotesize
\begin{center}
\renewcommand{\arraystretch}{1.7}
\begin{tabular}{ll}
differential equation
  & $(1-x^2)P_n'' - 2x\,P_n' + n(n+1)P_n = 0$ \\
series solution
  & $P_n(x)=\displaystyle\sum_{k=0}^{\lfloor n/2\rfloor}
       \frac{(-1)^k(2n-2k)!}{2^n k!\,(n-2k)!\,(n-k)!}\,x^{n-2k}$ \\
generating function
  & $\dfrac{1}{\sqrt{1-2xt+t^2}} = \displaystyle\sum_{n=0}^{\infty}P_n(x)\,t^n$ \\
recurrence
  & $(n+1)P_{n+1} = (2n+1)x\,P_n - n\,P_{n-1}$ \\
orthogonality
  & $\displaystyle\int_{-1}^{1}P_m(x)P_n(x)\,dx
       = \frac{2}{2n+1}\,\delta_{nm}$ \\
Rodrigues
  & $P_n(x)=\dfrac{1}{2^n n!}\dfrac{d^n}{dx^n}(x^2-1)^n$ \\
Schl\"afli
  & $P_n(x)=\dfrac{1}{2^{n+1}\pi i}
       \displaystyle\oint\frac{(t^2-1)^n}{(t-x)^{n+1}}\,dt$
\end{tabular}
\end{center}
\vspace{0.2em}
{\scriptsize The generating series converges for $|t|<1$ on
$x\in[-1,1]$.  In the Schl\"afli formula, the positively oriented contour
encloses $t=x$.}
\end{frame}

% ======================================================================
\section{Hermite Polynomial}
% ======================================================================

\begin{frame}[allowframebreaks]{The Hermite equation and series solution}
\footnotesize

The physicists' Hermite polynomials solve
\begin{equation}
  y''-2xy'+2ny=0,
  \qquad n=0,1,2,\ldots.
  \label{eq:hermite-eq}
\end{equation}
Let
\[
  y(x)=\sum_{j=0}^{\infty}a_jx^j.
\]
Termwise differentiation gives
\begin{align*}
  y'(x)
   &=\sum_{j=1}^{\infty}ja_jx^{j-1}
    =\sum_{j=0}^{\infty}(j+1)a_{j+1}x^j,\\
  y''(x)
   &=\sum_{j=2}^{\infty}j(j-1)a_jx^{j-2}
    =\sum_{j=0}^{\infty}(j+2)(j+1)a_{j+2}x^j.
\end{align*}
Also,
\[
  -2xy'=-2\sum_{j=1}^{\infty}ja_jx^j
         =-2\sum_{j=0}^{\infty}ja_jx^j,
\]
where the added $j=0$ term is zero, and
\[
  2ny=2n\sum_{j=0}^{\infty}a_jx^j.
\]
Substitution into \eqref{eq:hermite-eq} gives
\[
  \sum_{j=0}^{\infty}
  \left[(j+2)(j+1)a_{j+2}+2(n-j)a_j\right]x^j=0.
\]
Every coefficient must vanish, so
\begin{equation}
  (j+2)(j+1)a_{j+2}+2(n-j)a_j=0,
  \qquad
  a_{j+2}=\frac{2(j-n)}{(j+2)(j+1)}a_j.
  \label{eq:hermite-coefficient-recurrence}
\end{equation}
The recurrence links even coefficients only to even coefficients and odd coefficients only
to odd coefficients. For an integer $n\ge0$, the parity chain with the same parity as $n$
reaches $a_n$ and terminates because setting $j=n$ gives $a_{n+2}=0$. To obtain a
polynomial of degree exactly $n$, the initial coefficient of the opposite-parity chain is
set to zero.

We now derive every coefficient of the degree-$n$ polynomial. Set the conventional leading
coefficient to
\[
  a_n=2^n.
\]
In \eqref{eq:hermite-coefficient-recurrence}, put $j=n-2k-2$. Then
\begin{align*}
  a_{n-2k}
  &=\frac{2[(n-2k-2)-n]}{(n-2k)(n-2k-1)}a_{n-2k-2}\\
  &=-\frac{4(k+1)}{(n-2k)(n-2k-1)}a_{n-2k-2}.
\end{align*}
Solving for the lower-degree coefficient,
\begin{equation}
  a_{n-2k-2}
  =-\frac{(n-2k)(n-2k-1)}{4(k+1)}a_{n-2k}.
  \label{eq:hermite-downward}
\end{equation}
Applying \eqref{eq:hermite-downward} repeatedly gives
\begin{align*}
  a_{n-2k}
  &=(-1)^k
    \frac{n(n-1)(n-2)\cdots(n-2k+1)}{4^k k!}\,a_n\\
  &=(-1)^k\frac{n!}{(n-2k)!}\frac{2^n}{4^k k!}\\
  &=(-1)^k\frac{2^{n-2k}n!}{k!(n-2k)!}.
\end{align*}
Therefore
\begin{equation}
  H_n(x)=\sum_{k=0}^{\lfloor n/2\rfloor}
  (-1)^k(2x)^{n-2k}\frac{n!}{k!(n-2k)!}.
  \label{eq:hermite-series}
\end{equation}

For example, when $n=2$ the terms $k=0,1$ give
\begin{align*}
  H_2(x)
  &=\frac{(2x)^2\,2!}{0!\,2!}
    -\frac{(2x)^0\,2!}{1!\,0!}\\
  &=4x^2-2.
\end{align*}
When $n=3$ the terms $k=0,1$ give
\begin{align*}
  H_3(x)
  &=\frac{(2x)^3\,3!}{0!\,3!}
    -\frac{(2x)^1\,3!}{1!\,1!}\\
  &=8x^3-12x.
\end{align*}
Together with $H_0=1$ and $H_1=2x$, these agree with the stated normalization.

Multiplying \eqref{eq:hermite-eq} by $e^{-x^2}$ gives
\[
  e^{-x^2}y''-2xe^{-x^2}y'+2ne^{-x^2}y=0.
\]
Since
\[
  \frac{d}{dx}\!\left(e^{-x^2}y'\right)
  =e^{-x^2}y''-2xe^{-x^2}y',
\]
the Sturm--Liouville form is
\begin{equation}
  \frac{d}{dx}\!\left(e^{-x^2}y'\right)
  +2n e^{-x^2}y=0.
  \label{eq:hermite-sl}
\end{equation}

To prove orthogonality, write \eqref{eq:hermite-sl} for $H_n$ and $H_m$, multiply the
first by $H_m$, the second by $H_n$, and subtract:
\[
  H_m(e^{-x^2}H_n')'-H_n(e^{-x^2}H_m')'
  +2(n-m)e^{-x^2}H_nH_m=0.
\]
The first two terms form the derivative
\[
  \frac{d}{dx}\!\left[e^{-x^2}(H_mH_n'-H_nH_m')\right].
\]
Integrating over the real line gives
\begin{align*}
  &\left[e^{-x^2}(H_mH_n'-H_nH_m')\right]_{-\infty}^{\infty}\\
  &\qquad+2(n-m)\int_{-\infty}^{\infty}H_n(x)H_m(x)e^{-x^2}\,dx=0.
\end{align*}
The boundary term is zero because a polynomial times $e^{-x^2}$ tends to zero at both
infinities. Hence, for $n\ne m$,
\begin{equation}
  \int_{-\infty}^{\infty}H_n(x)H_m(x)e^{-x^2}\,dx=0.
  \label{eq:hermite-ortho}
\end{equation}
The generating-function calculation below supplies the diagonal norm.
\end{frame}

\begin{frame}{Shapes of the first Hermite polynomials}
\footnotesize
\begin{center}
  \includegraphics[width=0.86\linewidth]{figs/hermite-H.pdf}
\end{center}
The Hermite polynomials grow rapidly away from the origin. In the harmonic oscillator the
physical functions are $e^{-x^2/2}H_n(x)$, where the Gaussian factor restores decay.
\end{frame}

\begin{frame}[allowframebreaks]{Which equation does $\psi_n=e^{-x^2/2}H_n$ solve?}
\footnotesize

Let $y=H_n$ satisfy
\[
  y''-2xy'+2ny=0,
\]
and define
\[
  \psi_n(x)=e^{-x^2/2}y(x).
\]
Equivalently,
\[
  y(x)=e^{x^2/2}\psi_n(x).
\]
Using
\[
  \frac{d}{dx}e^{x^2/2}=xe^{x^2/2},
\]
the first derivative is
\begin{align*}
  y'
  &=xe^{x^2/2}\psi_n+e^{x^2/2}\psi_n'\\
  &=e^{x^2/2}(\psi_n'+x\psi_n).
\end{align*}
Differentiate again:
\begin{align*}
  y''
  &=xe^{x^2/2}(\psi_n'+x\psi_n)
    +e^{x^2/2}\frac{d}{dx}(\psi_n'+x\psi_n)\\
  &=e^{x^2/2}\left[x\psi_n'+x^2\psi_n
    +\psi_n''+\psi_n+x\psi_n'\right]\\
  &=e^{x^2/2}\left[\psi_n''+2x\psi_n'+(x^2+1)\psi_n\right].
\end{align*}
Substitute $y$, $y'$, and $y''$ into the Hermite equation and divide by the nonzero factor
$e^{x^2/2}$:
\begin{align*}
  0
  &=\psi_n''+2x\psi_n'+(x^2+1)\psi_n
    -2x(\psi_n'+x\psi_n)+2n\psi_n\\
  &=\psi_n''+2x\psi_n'+(x^2+1)\psi_n
    -2x\psi_n'-2x^2\psi_n+2n\psi_n\\
  &=\psi_n''+(2n+1-x^2)\psi_n.
\end{align*}
Therefore
\begin{equation}
  \psi_n''(x)+(2n+1-x^2)\psi_n(x)=0,
  \qquad
  \psi_n(x)=e^{-x^2/2}H_n(x).
  \label{eq:hermite-osc}
\end{equation}
\end{frame}

\begin{frame}[allowframebreaks]{Example: the quantum harmonic oscillator}
\footnotesize

For a particle of mass $M$ in a one-dimensional potential $V(x)$, the
time-dependent Schr\"odinger equation is
\begin{equation}
  i\hbar\frac{\partial\Psi}{\partial t}
   =
  \left[
    -\frac{\hbar^2}{2M}\frac{\partial^2}{\partial x^2}
    +V(x)
  \right]\Psi.
  \label{eq:schrod}
\end{equation}
Set
\[
  \Psi(x,t)=u(t)\psi(x).
\]
Then
\[
  \frac{\partial\Psi}{\partial t}=u'(t)\psi(x),
  \qquad
  \frac{\partial^2\Psi}{\partial x^2}=u(t)\psi''(x),
\]
so \eqref{eq:schrod} becomes
\[
  i\hbar u'(t)\psi(x)
   =
  u(t)\left[
    -\frac{\hbar^2}{2M}\psi''(x)+V(x)\psi(x)
  \right].
\]
For a nonzero separated solution, divide by $u(t)\psi(x)$:
\[
  i\hbar\frac{u'(t)}{u(t)}
   =
  \frac{-\dfrac{\hbar^2}{2M}\psi''(x)+V(x)\psi(x)}
       {\psi(x)}.
\]
The left side depends only on $t$ and the right side only on $x$, so both equal
a separation constant $E$.  The time equation is
\[
  i\hbar u'=Eu.
\]
Solving it step by step,
\begin{align*}
  \frac{u'}{u}
   &=-\frac{iE}{\hbar},\\
  \int\frac{u'}{u}\,dt
   &=-\frac{iE}{\hbar}\int dt,\\
  \log u
   &=-\frac{iE}{\hbar}t+C,\\
  u(t)
   &=Ae^{-iEt/\hbar},
  \qquad A=e^C.
\end{align*}
The spatial equation is
\begin{equation}
  -\frac{\hbar^2}{2M}\psi''(x)+V(x)\psi(x)=E\psi(x).
  \label{eq:time-independent-schrodinger}
\end{equation}

For the harmonic oscillator,
\[
  V(x)=\frac12Kx^2,
  \qquad
  \omega=\sqrt{\frac{K}{M}},
  \qquad
  K=M\omega^2,
\]
and therefore
\[
  V(x)=\frac12M\omega^2x^2.
\]
Substitution into \eqref{eq:time-independent-schrodinger} gives
\[
  -\frac{\hbar^2}{2M}\psi''
  +\frac12M\omega^2x^2\psi=E\psi.
\]
Multiply every term by $-2M/\hbar^2$:
\begin{align*}
  \left(-\frac{2M}{\hbar^2}\right)
  \left(-\frac{\hbar^2}{2M}\psi''\right)
   &+
  \left(-\frac{2M}{\hbar^2}\right)
  \left(\frac12M\omega^2x^2\psi\right)\\
   &=
  \left(-\frac{2M}{\hbar^2}\right)E\psi,
\end{align*}
so
\[
  \psi''
  -\frac{M^2\omega^2}{\hbar^2}x^2\psi
  =-\frac{2ME}{\hbar^2}\psi,
\]
or
\[
  \psi''
  +\left(
    \frac{2ME}{\hbar^2}
    -\frac{M^2\omega^2}{\hbar^2}x^2
  \right)\psi=0.
\]

Define the dimensionless coordinate
\[
  \xi=\sqrt{\frac{M\omega}{\hbar}}\,x.
\]
Because $\xi$ is linear in $x$,
\begin{align*}
  \frac{d\xi}{dx}
    &=\sqrt{\frac{M\omega}{\hbar}},\\
  \frac{d}{dx}
    &=\frac{d\xi}{dx}\frac{d}{d\xi}
      =\sqrt{\frac{M\omega}{\hbar}}\frac{d}{d\xi},\\
  \frac{d^2}{dx^2}
    &=\frac{M\omega}{\hbar}\frac{d^2}{d\xi^2}.
\end{align*}
Also,
\[
  x=\sqrt{\frac{\hbar}{M\omega}}\,\xi,
  \qquad
  x^2=\frac{\hbar}{M\omega}\xi^2,
\]
so
\[
  \frac{M^2\omega^2}{\hbar^2}x^2
   =
  \frac{M^2\omega^2}{\hbar^2}
  \frac{\hbar}{M\omega}\xi^2
   =
  \frac{M\omega}{\hbar}\xi^2.
\]
After substitution,
\[
  \frac{M\omega}{\hbar}\frac{d^2\psi}{d\xi^2}
  +\left(
    \frac{2ME}{\hbar^2}
    -\frac{M\omega}{\hbar}\xi^2
  \right)\psi=0.
\]
Divide by $M\omega/\hbar$:
\begin{align*}
  \frac{d^2\psi}{d\xi^2}
  +\left[
    \frac{2ME}{\hbar^2}\frac{\hbar}{M\omega}
    -\xi^2
  \right]\psi
  &=0,\\
  \frac{d^2\psi}{d\xi^2}
  +\left(
    \frac{2E}{\hbar\omega}-\xi^2
  \right)\psi
  &=0.
\end{align*}
Thus
\begin{equation}
  \frac{d^2\psi}{d\xi^2}
  +\left(\frac{2E}{\hbar\omega}-\xi^2\right)\psi=0.
  \label{eq:qho-dimensionless}
\end{equation}

For large $|\xi|$, the dominant equation is approximately
$\psi''-\xi^2\psi=0$, whose decaying behavior is Gaussian.  Factor it out by
setting
\[
  \psi(\xi)=e^{-\xi^2/2}y(\xi).
\]
The first derivative is
\begin{align*}
  \psi'
   &=
    \left(e^{-\xi^2/2}\right)'y
    +e^{-\xi^2/2}y'\\
   &=
    -\xi e^{-\xi^2/2}y
    +e^{-\xi^2/2}y'\\
   &=
    e^{-\xi^2/2}(y'-\xi y).
\end{align*}
Differentiate again:
\begin{align*}
  \psi''
   &=
    -\xi e^{-\xi^2/2}(y'-\xi y)
    +e^{-\xi^2/2}\frac{d}{d\xi}(y'-\xi y)\\
   &=
    e^{-\xi^2/2}
    \left[
      -\xi y'+\xi^2y+y''-y-\xi y'
    \right]\\
   &=
    e^{-\xi^2/2}
    \left[
      y''-2\xi y'+(\xi^2-1)y
    \right].
\end{align*}
Substitute these expressions into \eqref{eq:qho-dimensionless} and divide by
the nonzero factor $e^{-\xi^2/2}$:
\begin{align*}
  0
   &=
    y''-2\xi y'+(\xi^2-1)y
    +\left(\frac{2E}{\hbar\omega}-\xi^2\right)y\\
   &=
    y''-2\xi y'
    +\left(\frac{2E}{\hbar\omega}-1\right)y.
\end{align*}
Define
\[
  \varepsilon:=\frac{2E}{\hbar\omega}.
\]
Then
\begin{equation}
  y''-2\xi y'+(\varepsilon-1)y=0.
  \label{eq:qho-hermite-reduction}
\end{equation}

To determine when this equation has a polynomial solution, let
\[
  y(\xi)=\sum_{j=0}^{\infty}a_j\xi^j.
\]
Then
\begin{align*}
  y'
    &=\sum_{j=0}^{\infty}(j+1)a_{j+1}\xi^j,\\
  y''
    &=\sum_{j=0}^{\infty}(j+2)(j+1)a_{j+2}\xi^j,\\
  -2\xi y'
    &=-2\sum_{j=1}^{\infty}j\,a_j\xi^j
      =-2\sum_{j=0}^{\infty}j\,a_j\xi^j,\\
  (\varepsilon-1)y
    &=\sum_{j=0}^{\infty}(\varepsilon-1)a_j\xi^j.
\end{align*}
Hence
\[
  \sum_{j=0}^{\infty}
  \left[
    (j+2)(j+1)a_{j+2}
    +(\varepsilon-1-2j)a_j
  \right]\xi^j=0,
\]
so every coefficient satisfies
\begin{equation}
  a_{j+2}
    =\frac{2j+1-\varepsilon}{(j+2)(j+1)}a_j.
  \label{eq:qho-series-recurrence}
\end{equation}

The even and odd coefficient chains are independent.  If a chain does not
terminate, then for large $j$ of the appropriate parity,
\begin{align*}
  \frac{a_{j+2}}{a_j}
   &=
    \frac{2j+1-\varepsilon}{(j+2)(j+1)}
    \sim\frac{2j}{j^2}
    =\frac{2}{j}.
\end{align*}
For comparison,
\[
  e^{\xi^2}
    =\sum_{k=0}^{\infty}\frac{\xi^{2k}}{k!}
\]
has successive even-power coefficient ratio
\[
  \frac{1/(k+1)!}{1/k!}
   =\frac{1}{k+1}
   =\frac{2}{2k+2}
   \sim\frac{2}{j}
   \qquad (j=2k),
\]
and $\xi e^{\xi^2}$ has the analogous odd-power ratio.  Thus a
nonterminating parity solution has the growing asymptotic component
$y(\xi)\sim e^{\xi^2}$ (up to an algebraic factor).  Consequently
\[
  \psi(\xi)=e^{-\xi^2/2}y(\xi)
\]
contains an $e^{+\xi^2/2}$ component and is not square-integrable on the real
line.  A bound-state parity solution must therefore terminate.

Termination at degree $n$ occurs when the numerator of
\eqref{eq:qho-series-recurrence} vanishes at $j=n$:
\begin{align*}
  2n+1-\varepsilon&=0,\\
  \varepsilon&=2n+1.
\end{align*}
For this value,
\[
  y''-2\xi y'+2ny=0,
\]
which is the physicists' Hermite equation.  Selecting the parity chain that
contains $\xi^n$ and setting the other initial coefficient to zero gives
$y(\xi)=H_n(\xi)$ up to normalization.  Since
$\varepsilon=2E/(\hbar\omega)$,
\begin{align*}
  \frac{2E_n}{\hbar\omega}
    &=2n+1,\\
  E_n
    &=\frac{\hbar\omega}{2}(2n+1)
      =\left(n+\frac12\right)\hbar\omega.
\end{align*}
Therefore
\begin{equation}
  E_n=\left(n+\frac12\right)\hbar\omega,
  \qquad n=0,1,2,\ldots .
  \label{eq:energy-levels}
\end{equation}

Write
\[
  \psi_n(x)=N_n e^{-\xi^2/2}H_n(\xi),
  \qquad
  \xi=\sqrt{\frac{M\omega}{\hbar}}\,x.
\]
Because
\[
  dx=\sqrt{\frac{\hbar}{M\omega}}\,d\xi
\]
and
\[
  \int_{-\infty}^{\infty}
    e^{-\xi^2}H_n(\xi)^2\,d\xi
    =2^n n!\sqrt{\pi},
\]
normalization requires
\begin{align*}
  1
   &=\int_{-\infty}^{\infty}|\psi_n(x)|^2\,dx\\
   &=|N_n|^2
     \sqrt{\frac{\hbar}{M\omega}}
     \int_{-\infty}^{\infty}
       e^{-\xi^2}H_n(\xi)^2\,d\xi\\
   &=|N_n|^2
     \sqrt{\frac{\hbar}{M\omega}}
     2^n n!\sqrt{\pi}.
\end{align*}
Choosing $N_n$ real and positive,
\begin{align*}
  N_n^2
   &=
    \sqrt{\frac{M\omega}{\hbar}}
    \frac{1}{2^n n!\sqrt{\pi}},\\
  N_n
   &=
    \left(\frac{M\omega}{\pi\hbar}\right)^{1/4}
    \frac{1}{\sqrt{2^n n!}}.
\end{align*}
Thus the normalized spatial eigenfunction is
\begin{equation}
  \psi_n(x)
    =
  \left(\frac{M\omega}{\pi\hbar}\right)^{1/4}
  \frac{e^{-\xi^2/2}H_n(\xi)}{\sqrt{2^n n!}},
  \qquad
  \xi=\sqrt{\frac{M\omega}{\hbar}}\,x.
  \label{eq:qho-normalized-eigenfunction}
\end{equation}
A normalized state expanded in this orthonormal basis is
\begin{equation}
  \Psi(x,t)
    =\sum_{n=0}^{\infty}
      c_n\psi_n(x)e^{-iE_nt/\hbar},
  \qquad
  \sum_{n=0}^{\infty}|c_n|^2=1.
  \label{eq:qho-solution}
\end{equation}
\end{frame}

\begin{frame}[allowframebreaks]{Generating function, norm, Rodrigues, and Schl\"afli}
\footnotesize

Starting from the explicit Hermite series,
\[
  H_n(x)=\sum_{k=0}^{\lfloor n/2\rfloor}
  (-1)^k(2x)^{n-2k}\frac{n!}{k!(n-2k)!},
\]
form the exponential generating function:
\[
  \sum_{n=0}^{\infty}H_n(x)\frac{t^n}{n!}
  =\sum_{n=0}^{\infty}\sum_{k=0}^{\lfloor n/2\rfloor}
    (-1)^k\frac{(2x)^{n-2k}t^n}{k!(n-2k)!}.
\]
Set $j=n-2k$, so $n=j+2k$ and $j,k\ge0$. Then
\begin{align*}
  \sum_{n=0}^{\infty}H_n(x)\frac{t^n}{n!}
  &=\sum_{k=0}^{\infty}\sum_{j=0}^{\infty}
    (-1)^k\frac{(2x)^j t^{j+2k}}{k!j!}\\
  &=\left(\sum_{k=0}^{\infty}\frac{(-1)^k t^{2k}}{k!}\right)
    \left(\sum_{j=0}^{\infty}\frac{(2xt)^j}{j!}\right)\\
  &=e^{-t^2}e^{2xt}.
\end{align*}
Thus
\begin{equation}
  g(x,t)=e^{-t^2+2xt}
  =\sum_{n=0}^{\infty}H_n(x)\frac{t^n}{n!}.
  \label{eq:hermite-gen}
\end{equation}

We now compute the norm.  Take real $s$ and $t$ near the origin, multiply two
generating functions, and integrate with the weight:
\begin{align*}
  I(s,t)
  &:=\int_{-\infty}^{\infty}e^{-x^2}g(x,s)g(x,t)\,dx\\
  &=\int_{-\infty}^{\infty}
    e^{-x^2-s^2+2sx-t^2+2tx}\,dx.
\end{align*}
Complete the square:
\begin{align*}
  -x^2+2(s+t)x-s^2-t^2
  &=-\bigl[x-(s+t)\bigr]^2+(s+t)^2-s^2-t^2\\
  &=-\bigl[x-(s+t)\bigr]^2+2st.
\end{align*}
Therefore
\begin{align*}
  I(s,t)
  &=e^{2st}\int_{-\infty}^{\infty}e^{-[x-(s+t)]^2}\,dx\\
  &=\sqrt\pi\,e^{2st}\\
  &=\sqrt\pi\sum_{k=0}^{\infty}\frac{2^k s^k t^k}{k!}.
\end{align*}
On the other hand, expanding the two generating functions first gives
\[
  I(s,t)
  =\sum_{m=0}^{\infty}\sum_{n=0}^{\infty}
    \frac{s^m t^n}{m!n!}
    \int_{-\infty}^{\infty}H_m(x)H_n(x)e^{-x^2}\,dx.
\]
For $m\ne n$, the coefficient of $s^mt^n$ on $\sqrt\pi e^{2st}$ is zero. For $m=n$,
coefficient comparison gives
\[
  \frac{1}{(n!)^2}
  \int_{-\infty}^{\infty}H_n(x)^2e^{-x^2}\,dx
  =\sqrt\pi\frac{2^n}{n!}.
\]
Hence
\begin{equation}
  \int_{-\infty}^{\infty}H_m(x)H_n(x)e^{-x^2}\,dx
  =2^n n!\sqrt\pi\,\delta_{mn}.
  \label{eq:hermite-norm}
\end{equation}
Equivalently, the dimensionless oscillator functions
\[
  \widehat\psi_n(x)
  =\frac{e^{-x^2/2}H_n(x)}{\pi^{1/4}\sqrt{2^n n!}}
\]
satisfy $\int_{-\infty}^{\infty}\widehat\psi_m\widehat\psi_n\,dx=\delta_{mn}$.

To derive Rodrigues' formula, rewrite the generating function as
\[
  g(x,t)=e^{x^2}e^{-(t-x)^2}.
\]
Differentiate $m$ times with respect to $t$:
\begin{align*}
  \frac{\partial^m g}{\partial t^m}
  &=e^{x^2}\frac{d^m}{dt^m}e^{-(t-x)^2},\\
  \frac{\partial^m}{\partial t^m}
  \left(\sum_{n=0}^{\infty}H_n(x)\frac{t^n}{n!}\right)
  &=\sum_{n=m}^{\infty}H_n(x)\frac{t^{n-m}}{(n-m)!}\\
  &=H_m(x)+H_{m+1}(x)t+\frac{H_{m+2}(x)}{2!}t^2+\cdots.
\end{align*}
Setting $t=0$ leaves
\[
  H_m(x)=e^{x^2}\left.\frac{d^m}{dt^m}e^{-(t-x)^2}\right|_{t=0}.
\]
Let $s=x-t$. Then $ds/dt=-1$, so
\[
  \frac{d^m}{dt^m}=(-1)^m\frac{d^m}{ds^m},
  \qquad t=0\Longrightarrow s=x.
\]
Thus
\begin{equation}
  H_m(x)=(-1)^m e^{x^2}\frac{d^m}{dx^m}e^{-x^2}.
  \label{eq:hermite-rodrigues}
\end{equation}

Finally, divide \eqref{eq:hermite-gen} by $t^{m+1}$ and integrate around a positively
oriented contour $\gamma$ enclosing the origin:
\begin{align*}
  \oint_\gamma\frac{e^{-t^2+2xt}}{t^{m+1}}\,dt
  &=\sum_{n=0}^{\infty}\frac{H_n(x)}{n!}
    \oint_\gamma t^{n-m-1}\,dt\\
  &=\frac{H_m(x)}{m!}(2\pi i).
\end{align*}
Therefore
\begin{equation}
  H_m(x)=\frac{m!}{2\pi i}
  \oint_\gamma\frac{e^{-t^2+2xt}}{t^{m+1}}\,dt.
  \label{eq:hermite-schlafli}
\end{equation}
Because the exponential is entire, any sufficiently small positively oriented contour
around $t=0$ is admissible.
\end{frame}

\begin{frame}[allowframebreaks]{Example: special values $H_n(0)$}
\footnotesize

Set $x=0$ in the generating function:
\[
  e^{-t^2}
  =\sum_{n=0}^{\infty}H_n(0)\frac{t^n}{n!}.
\]
The ordinary exponential series gives
\[
  e^{-t^2}
  =\sum_{k=0}^{\infty}\frac{(-t^2)^k}{k!}
  =\sum_{k=0}^{\infty}(-1)^k\frac{t^{2k}}{k!}.
\]
There are no odd powers on the right, so the coefficient of $t^{2k+1}$ satisfies
\[
  \frac{H_{2k+1}(0)}{(2k+1)!}=0,
  \qquad
  H_{2k+1}(0)=0.
\]
For the even powers,
\[
  \frac{H_{2k}(0)}{(2k)!}=\frac{(-1)^k}{k!},
\]
and hence
\begin{equation}
  H_{2k+1}(0)=0,
  \qquad
  H_{2k}(0)=(-1)^k\frac{(2k)!}{k!}.
  \label{eq:hermite-zero}
\end{equation}
Checks from the explicit polynomials are
\[
  H_0(0)=1,
  \qquad H_1(0)=0,
  \qquad H_2(0)=-2,
  \qquad H_4(0)=12,
\]
which agree with \eqref{eq:hermite-zero}.
\end{frame}

\begin{frame}[allowframebreaks]{Recurrence relations}
\footnotesize

Start from
\[
  g(x,t)=e^{-t^2+2xt}
  =\sum_{n=0}^{\infty}H_n(x)\frac{t^n}{n!}.
\]
Differentiate with respect to $x$. Direct differentiation gives
\[
  \frac{\partial g}{\partial x}=2t\,g.
\]
Termwise differentiation gives
\[
  \frac{\partial g}{\partial x}
  =\sum_{n=0}^{\infty}H_n'(x)\frac{t^n}{n!}.
\]
Also,
\begin{align*}
  2t\,g
  &=2\sum_{j=0}^{\infty}H_j(x)\frac{t^{j+1}}{j!}\\
  &=2\sum_{n=1}^{\infty}H_{n-1}(x)\frac{t^n}{(n-1)!}\\
  &=2\sum_{n=1}^{\infty}nH_{n-1}(x)\frac{t^n}{n!}.
\end{align*}
Matching coefficients of $t^n/n!$ gives $H_0'=0$ and, for $n\ge1$,
\begin{equation}
  H_n'(x)=2nH_{n-1}(x).
  \label{eq:hermite-deriv}
\end{equation}

Now differentiate the generating function with respect to $t$. From the exponential,
\[
  \frac{\partial g}{\partial t}=(2x-2t)g.
\]
From the series,
\begin{align*}
  \frac{\partial g}{\partial t}
  &=\sum_{j=1}^{\infty}H_j(x)\frac{j t^{j-1}}{j!}\\
  &=\sum_{n=0}^{\infty}H_{n+1}(x)\frac{t^n}{n!}.
\end{align*}
Expand the right side of $(2x-2t)g$:
\begin{align*}
  (2x-2t)g
  &=2x\sum_{n=0}^{\infty}H_n\frac{t^n}{n!}
    -2\sum_{j=0}^{\infty}H_j\frac{t^{j+1}}{j!}\\
  &=2x\sum_{n=0}^{\infty}H_n\frac{t^n}{n!}
    -2\sum_{n=1}^{\infty}nH_{n-1}\frac{t^n}{n!}.
\end{align*}
Matching coefficients of $t^n/n!$ yields
\begin{equation}
  H_{n+1}(x)=2xH_n(x)-2nH_{n-1}(x),
  \qquad n\ge0,
  \label{eq:hermite-recur}
\end{equation}
where the $n=0$ case is $H_1=2xH_0=2x$ and does not require $H_{-1}$.
\end{frame}

\begin{frame}{Useful formulas for the Hermite polynomial}
\footnotesize
\begin{center}
\renewcommand{\arraystretch}{1.7}
\begin{tabular}{ll}
differential equation
  & $H_n'' - 2x\,H_n' + 2n\,H_n = 0$ \\
series solution
  & $H_n(x)=\displaystyle\sum_{k=0}^{\lfloor n/2\rfloor}
       (-1)^k(2x)^{n-2k}\,\frac{n!}{k!\,(n-2k)!}$ \\
generating function
  & $e^{-t^2+2tx} = \displaystyle\sum_{n=0}^{\infty}H_n(x)\,\frac{t^n}{n!}$ \\
recurrence
  & $H_n' = 2n\,H_{n-1}$ \\
  & $H_{n+1}(x) = 2x\,H_n(x) - 2n\,H_{n-1}(x)$ \\
orthogonality
  & $\displaystyle\int_{-\infty}^{\infty}H_mH_n\,e^{-x^2}dx
       = n!\,2^n\sqrt{\pi}\,\delta_{mn}$ \\
Rodrigues
  & $H_n(x)=(-1)^ne^{x^2}\dfrac{d^n}{dx^n}e^{-x^2}$ \\
Schl\"afli
  & $H_n(x)=\dfrac{n!}{2\pi i}
       \displaystyle\oint\frac{e^{-t^2+2tx}}{t^{n+1}}\,dt$
\end{tabular}
\end{center}
\vspace{0.2em}
{\scriptsize The Hermite generating function is entire in $t$.  In the
contour formula, the positively oriented contour encloses $t=0$.}
\end{frame}

% ======================================================================
\section{Laguerre Polynomial}
% ======================================================================

\begin{frame}[allowframebreaks]{The Laguerre equation and series solution}
\footnotesize

The standard Laguerre equation is
\begin{equation}
  x\,y''(x)+(1-x)y'(x)+n\,y(x)=0,
  \qquad n=0,1,2,\ldots .
  \label{eq:laguerre-eq}
\end{equation}
Seek a solution analytic at $x=0$ in the form
\[
  y(x)=\sum_{j=0}^{\infty}a_jx^j .
\]
Termwise differentiation gives
\begin{align*}
  y'(x)
    &=\sum_{j=1}^{\infty}j\,a_jx^{j-1}
      =\sum_{k=0}^{\infty}(k+1)a_{k+1}x^k,\\
  y''(x)
    &=\sum_{j=2}^{\infty}j(j-1)a_jx^{j-2}.
\end{align*}
Multiplication by $x$ and the index shift $k=j-1$ give
\begin{align*}
  x\,y''(x)
    &=\sum_{j=2}^{\infty}j(j-1)a_jx^{j-1}\\
    &=\sum_{k=1}^{\infty}(k+1)k\,a_{k+1}x^k\\
    &=\sum_{k=0}^{\infty}(k+1)k\,a_{k+1}x^k,
\end{align*}
where the inserted $k=0$ term is zero.  Similarly,
\[
  -x\,y'(x)
   =-\sum_{j=1}^{\infty}j\,a_jx^j
   =-\sum_{k=0}^{\infty}k\,a_kx^k,
  \qquad
  n\,y(x)=\sum_{k=0}^{\infty}n\,a_kx^k .
\]
Substitution into \eqref{eq:laguerre-eq} therefore yields
\begin{align*}
  0
   &=\sum_{k=0}^{\infty}
      \left[
        (k+1)k\,a_{k+1}
        +(k+1)a_{k+1}
        -k\,a_k+n\,a_k
      \right]x^k\\
   &=\sum_{k=0}^{\infty}
      \left[
        (k+1)^2a_{k+1}+(n-k)a_k
      \right]x^k .
\end{align*}
A power series is identically zero only when every coefficient is zero, so
\begin{equation}
  (k+1)^2a_{k+1}+(n-k)a_k=0,
  \qquad
  a_{k+1}=-\frac{n-k}{(k+1)^2}\,a_k .
  \label{eq:laguerre-coeff-recurrence}
\end{equation}

Choose the standard normalization $a_0=1$.  Iterating the recurrence gives, for
$1\le k\le n$,
\begin{align*}
  a_k
   &=(-1)^k
     \frac{n(n-1)\cdots(n-k+1)}
          {1^2\,2^2\cdots k^2}\,a_0\\
   &=(-1)^k
     \frac{n!/(n-k)!}{(k!)^2}\\
   &=\frac{(-1)^k n!}{(n-k)!(k!)^2}.
\end{align*}
At $k=n$, \eqref{eq:laguerre-coeff-recurrence} gives
\[
  a_{n+1}=-\frac{n-n}{(n+1)^2}a_n=0.
\]
All later coefficients are then zero, because the recurrence successively gives
$a_{n+2}=a_{n+3}=\cdots=0$.  Thus the analytic solution terminates:
\begin{equation}
  L_n(x)
   =\sum_{k=0}^{n}
      \frac{(-1)^k n!}{(n-k)!(k!)^2}\,x^k
   =\sum_{k=0}^{n}
      (-1)^k\binom{n}{k}\frac{x^k}{k!}.
  \label{eq:laguerre-series}
\end{equation}
For example,
\begin{align*}
  L_0(x)
    &=1,\\
  L_1(x)
    &=1-x,\\
  L_2(x)
    &=1-2x+\frac{x^2}{2},\\
  L_3(x)
    &=1-3x+\frac{3}{2}x^2-\frac{1}{6}x^3.
\end{align*}
As a direct check for $L_2$,
\begin{align*}
  L_2'(x)&=-2+x,
  &
  L_2''(x)&=1,\\
  xL_2''+(1-x)L_2'+2L_2
    &=x+(1-x)(-2+x)+2\left(1-2x+\frac{x^2}{2}\right)\\
    &=x+\left(-2+3x-x^2\right)
      +\left(2-4x+x^2\right)\\
    &=0.
\end{align*}

The generalized Laguerre equation is
\[
  x\,y''+(\alpha+1-x)y'+n\,y=0.
\]
Its polynomial solution is denoted by $L_n^{(\alpha)}(x)$.  The standard
orthogonality theory on $[0,\infty)$ requires $\alpha>-1$; the present chapter uses
the case $\alpha=0$, for which $L_n^{(0)}=L_n$.

\medskip
To expose the Sturm--Liouville structure of the standard equation, multiply
\eqref{eq:laguerre-eq} by $e^{-x}$.  Since
\[
  \frac{d}{dx}\!\left(xe^{-x}\right)
    =e^{-x}-xe^{-x}
    =(1-x)e^{-x},
\]
the product rule gives
\[
  \frac{d}{dx}\!\left(xe^{-x}y'\right)
    =xe^{-x}y''+(1-x)e^{-x}y'.
\]
Hence
\begin{equation}
  \frac{d}{dx}\!\left(xe^{-x}y'\right)
   +n e^{-x}y=0 .
  \label{eq:laguerre-sl}
\end{equation}

Let $L_n$ and $L_m$ satisfy \eqref{eq:laguerre-sl} with eigenvalues $n$ and
$m$.  Their two equations are
\begin{align*}
  \left(xe^{-x}L_n'\right)'+n e^{-x}L_n&=0,\\
  \left(xe^{-x}L_m'\right)'+m e^{-x}L_m&=0.
\end{align*}
Multiply the first by $L_m$, the second by $L_n$, and subtract:
\[
  L_m\left(xe^{-x}L_n'\right)'
  -L_n\left(xe^{-x}L_m'\right)'
  +(n-m)e^{-x}L_nL_m=0.
\]
The first two terms are one total derivative, because
\begin{align*}
  \frac{d}{dx}\!
    \left[xe^{-x}\left(L_mL_n'-L_nL_m'\right)\right]
  &=
    L_m\left(xe^{-x}L_n'\right)'
    +xe^{-x}L_m'L_n'\\
  &\quad
    -L_n\left(xe^{-x}L_m'\right)'
    -xe^{-x}L_n'L_m'\\
  &=
    L_m\left(xe^{-x}L_n'\right)'
    -L_n\left(xe^{-x}L_m'\right)'.
\end{align*}
Integrating from $0$ to $R$ gives
\[
  \left[
    xe^{-x}\left(L_mL_n'-L_nL_m'\right)
  \right]_{0}^{R}
  +(n-m)\int_0^R e^{-x}L_n(x)L_m(x)\,dx=0.
\]
At $x=0$, the boundary expression vanishes because of the factor $x$.  As
$R\to\infty$, each $L_j$ and $L_j'$ is a polynomial, while $e^{-R}$ decays faster
than any power of $R$; consequently
\[
  \lim_{R\to\infty}
  R\,e^{-R}\left(L_m(R)L_n'(R)-L_n(R)L_m'(R)\right)=0.
\]
Taking $R\to\infty$ therefore gives
\[
  (n-m)\int_0^\infty e^{-x}L_n(x)L_m(x)\,dx=0.
\]
Thus, when $n\ne m$,
\begin{equation}
  \int_0^\infty L_n(x)L_m(x)e^{-x}\,dx=0.
  \label{eq:laguerre-ortho}
\end{equation}
The diagonal norm is derived from the generating function below.
\end{frame}

\begin{frame}{Shapes of the first Laguerre polynomials}
\footnotesize
\begin{center}
  \includegraphics[width=0.86\linewidth]{figs/laguerre-L.pdf}
\end{center}
All standard Laguerre polynomials satisfy $L_n(0)=1$. Their positive zeros and the weight
$e^{-x}$ are what make them natural on the half-line $[0,\infty)$.
\end{frame}

\begin{frame}[allowframebreaks]{Generating function, norm, Schl\"afli, and Rodrigues}
\footnotesize

Start from the finite series \eqref{eq:laguerre-series}:
\[
  L_n(x)=\sum_{k=0}^{n}(-1)^k\binom{n}{k}\frac{x^k}{k!}.
\]
For $|t|<1$, sum over $n$ and interchange the sums:
\begin{align*}
  \sum_{n=0}^{\infty}L_n(x)t^n
   &=
     \sum_{n=0}^{\infty}
     \sum_{k=0}^{n}
       (-1)^k\binom{n}{k}\frac{x^k}{k!}t^n\\
   &=
     \sum_{k=0}^{\infty}\frac{(-x)^k}{k!}
     \sum_{n=k}^{\infty}\binom{n}{k}t^n.
\end{align*}
The inner sum follows by differentiating the geometric series $k$ times:
\begin{align*}
  \frac{1}{1-t}
    &=\sum_{n=0}^{\infty}t^n,\\
  \frac{d^k}{dt^k}\frac{1}{1-t}
    &=\frac{k!}{(1-t)^{k+1}}
      =\sum_{n=k}^{\infty}\frac{n!}{(n-k)!}t^{n-k}.
\end{align*}
Multiplying the last equality by $t^k/k!$ gives
\[
  \sum_{n=k}^{\infty}\binom{n}{k}t^n
    =\frac{t^k}{(1-t)^{k+1}}.
\]
Substitution into the double sum now yields
\begin{align*}
  \sum_{n=0}^{\infty}L_n(x)t^n
    &=
      \sum_{k=0}^{\infty}
      \frac{(-x)^k}{k!}
      \frac{t^k}{(1-t)^{k+1}}\\
    &=
      \frac{1}{1-t}
      \sum_{k=0}^{\infty}
      \frac{1}{k!}
      \left(-\frac{xt}{1-t}\right)^k\\
    &=
      \frac{1}{1-t}
      \exp\!\left(-\frac{xt}{1-t}\right).
\end{align*}
Therefore
\begin{equation}
  g(x,t)
    :=\frac{1}{1-t}\exp\!\left(-\frac{xt}{1-t}\right)
    =\sum_{n=0}^{\infty}L_n(x)t^n,
  \qquad |t|<1.
  \label{eq:laguerre-gen}
\end{equation}
Setting $x=0$ gives
\[
  \sum_{n=0}^{\infty}L_n(0)t^n
    =g(0,t)
    =\frac{1}{1-t}
    =\sum_{n=0}^{\infty}t^n,
\]
so coefficient comparison confirms $L_n(0)=1$.

\medskip
Cauchy's coefficient formula applied to
$g(x,z)=\sum_{n=0}^{\infty}L_n(x)z^n$ gives
\begin{align*}
  L_m(x)
    &=[z^m]\,g(x,z)\\
    &=\frac{1}{2\pi i}
      \oint_{|z|=\rho}\frac{g(x,z)}{z^{m+1}}\,dz,
      \qquad 0<\rho<1.
\end{align*}
After inserting $g$, this becomes the Schl\"afli contour representation
\begin{equation}
  L_m(x)
    =\frac{1}{2\pi i}
      \oint_{|z|=\rho}
      \frac{\exp\!\left(-xz/(1-z)\right)}
           {(1-z)z^{m+1}}\,dz,
  \qquad 0<\rho<1.
  \label{eq:laguerre-schlafli}
\end{equation}
Indeed, expanding the integrand verifies the coefficient extraction directly:
\[
  \frac{g(x,z)}{z^{m+1}}
    =\sum_{n=0}^{\infty}L_n(x)z^{n-m-1},
\]
and only the term with $n=m$ has residue $L_m(x)$ at the origin.

\medskip
The generating function also determines the diagonal norm.  Take real $s$ and
$t$ in a sufficiently small neighborhood of the origin, so all integrals below
converge absolutely, and define
\[
  I(s,t)
    :=\int_0^\infty e^{-x}g(x,s)g(x,t)\,dx.
\]
Substituting \eqref{eq:laguerre-gen} gives
\begin{align*}
  I(s,t)
   &=
    \frac{1}{(1-s)(1-t)}
    \int_0^\infty
      \exp\!\left[
        -x-\frac{xs}{1-s}-\frac{xt}{1-t}
      \right]dx\\
   &=
    \frac{1}{(1-s)(1-t)}
    \int_0^\infty
      \exp\!\left[
        -x\left(
          1+\frac{s}{1-s}+\frac{t}{1-t}
        \right)
      \right]dx.
\end{align*}
The coefficient in the exponential simplifies as
\begin{align*}
  1+\frac{s}{1-s}+\frac{t}{1-t}
   &=
    \frac{(1-s)(1-t)+s(1-t)+t(1-s)}
         {(1-s)(1-t)}\\
   &=
    \frac{1-s-t+st+s-st+t-st}
         {(1-s)(1-t)}\\
   &=
    \frac{1-st}{(1-s)(1-t)}.
\end{align*}
Using $\int_0^\infty e^{-Ax}\,dx=1/A$ when $\operatorname{Re}A>0$,
\begin{align*}
  I(s,t)
   &=
    \frac{1}{(1-s)(1-t)}
    \frac{(1-s)(1-t)}{1-st}\\
   &=\frac{1}{1-st}
    =\sum_{r=0}^{\infty}(st)^r.
\end{align*}
On the other hand, expanding both generating functions before integration gives
\begin{align*}
  I(s,t)
   &=
    \int_0^\infty e^{-x}
      \left(\sum_{m=0}^{\infty}L_m(x)s^m\right)
      \left(\sum_{n=0}^{\infty}L_n(x)t^n\right)dx\\
   &=
    \sum_{m=0}^{\infty}\sum_{n=0}^{\infty}
      \left[
        \int_0^\infty e^{-x}L_m(x)L_n(x)\,dx
      \right]s^m t^n.
\end{align*}
Since
\[
  \frac{1}{1-st}
   =\sum_{r=0}^{\infty}s^rt^r
   =\sum_{m=0}^{\infty}\sum_{n=0}^{\infty}
      \delta_{mn}s^mt^n,
\]
comparison of the coefficient of $s^mt^n$ yields the complete orthonormality
formula
\begin{equation}
  \int_0^\infty e^{-x}L_m(x)L_n(x)\,dx=\delta_{mn}.
  \label{eq:laguerre-norm}
\end{equation}

\medskip
Finally, derive Rodrigues' formula from
\eqref{eq:laguerre-schlafli}.  Set
\[
  z=\frac{s-x}{s}=1-\frac{x}{s}.
\]
Then
\[
  1-z=\frac{x}{s},
  \qquad
  z^{m+1}=\frac{(s-x)^{m+1}}{s^{m+1}},
  \qquad
  dz=\frac{x}{s^2}\,ds,
\]
and
\[
  -\frac{xz}{1-z}
   =-x
     \frac{(s-x)/s}{x/s}
   =-(s-x)
   =x-s.
\]
Consequently,
\begin{align*}
  \frac{\exp\!\left(-xz/(1-z)\right)}
       {(1-z)z^{m+1}}\,dz
   &=
    \frac{e^{x-s}}
         {(x/s)\bigl((s-x)/s\bigr)^{m+1}}
    \frac{x}{s^2}\,ds\\
   &=
    e^x\,
    \frac{s^m e^{-s}}{(s-x)^{m+1}}\,ds.
\end{align*}
A small positively oriented circle around $z=0$ maps to a small positively
oriented contour around $s=x$.  Thus
\begin{align*}
  L_m(x)
   &=
    \frac{e^x}{2\pi i}
    \oint
      \frac{s^m e^{-s}}{(s-x)^{m+1}}\,ds\\
   &=
    \frac{e^x}{m!}
    \left.
      \frac{d^m}{ds^m}\left(s^m e^{-s}\right)
    \right|_{s=x},
\end{align*}
where the last step is Cauchy's differentiation formula.  Therefore
\begin{equation}
  L_m(x)
    =\frac{e^x}{m!}
      \frac{d^m}{dx^m}\left(x^m e^{-x}\right).
  \label{eq:laguerre-rodrigues}
\end{equation}
The change of variables used $x\ne0$; both sides are polynomials, so the identity
also holds at $x=0$ by continuity.
\end{frame}

\begin{frame}[allowframebreaks]{Recurrence relations for Laguerre polynomials}
\footnotesize

Differentiate the generating function \eqref{eq:laguerre-gen} with respect to $x$:
\begin{align*}
  \frac{\partial g}{\partial x}
   &=
    \frac{1}{1-t}
    \exp\!\left(-\frac{xt}{1-t}\right)
    \left(-\frac{t}{1-t}\right)\\
   &=
    -\frac{t}{1-t}\,g(x,t).
\end{align*}
Using
\[
  g(x,t)=\sum_{n=0}^{\infty}L_n(x)t^n,
  \qquad
  \frac{\partial g}{\partial x}
    =\sum_{n=0}^{\infty}L_n'(x)t^n,
\]
we obtain
\[
  (1-t)\sum_{n=0}^{\infty}L_n'(x)t^n
    =-t\sum_{n=0}^{\infty}L_n(x)t^n.
\]
Expand the left-hand side and shift indices:
\begin{align*}
  \sum_{n=0}^{\infty}L_n't^n
   -\sum_{n=0}^{\infty}L_n't^{n+1}
   &=
   -\sum_{n=0}^{\infty}L_nt^{n+1},\\
  L_0'
   +\sum_{n=1}^{\infty}L_n't^n
   -\sum_{n=1}^{\infty}L_{n-1}'t^n
   &=
   -\sum_{n=1}^{\infty}L_{n-1}t^n.
\end{align*}
The coefficient of $t^0$ gives $L_0'=0$.  For every $n\ge1$, the coefficient
of $t^n$ gives
\[
  L_n'-L_{n-1}'=-L_{n-1},
\]
or
\begin{equation}
  L_n'(x)=L_{n-1}'(x)-L_{n-1}(x),
  \qquad n\ge1.
  \label{eq:laguerre-derivative-shift}
\end{equation}

A second derivative identity follows directly from the finite series.  From
\eqref{eq:laguerre-series},
\begin{equation}
\begin{aligned}
  xL_n'(x)
   &=
    x\sum_{k=1}^{n}
      \frac{(-1)^k n!}{(n-k)!(k!)^2}
      kx^{k-1}\\
   &=
    \sum_{k=1}^{n}
      \frac{(-1)^k n!k}{(n-k)!(k!)^2}x^k.
\end{aligned}
\label{eq:laguerre-x-derivative-expansion}
\end{equation}
Now compare this with $n(L_n-L_{n-1})$.  For $k=0$, the coefficient is
$n(1-1)=0$.  For $1\le k\le n-1$, the coefficient of $x^k$ is
\begin{align*}
  &n(-1)^k
  \left[
    \frac{n!}{(n-k)!(k!)^2}
    -\frac{(n-1)!}{(n-1-k)!(k!)^2}
  \right]\\
  &\quad=
  \frac{n(-1)^k(n-1)!}{(n-k)!(k!)^2}
  \left[n-(n-k)\right]\\
  &\quad=
  \frac{(-1)^k n!k}{(n-k)!(k!)^2}.
\end{align*}
For $k=n$, $L_{n-1}$ has no $x^n$ term, and the coefficient is
\begin{align*}
  n\left[
    \frac{(-1)^n n!}{0!(n!)^2}
  \right]
   &=
    \frac{(-1)^n n!n}{0!(n!)^2},
\end{align*}
which is again the $k=n$ coefficient in
\eqref{eq:laguerre-x-derivative-expansion}.  Thus all coefficients agree, and
\begin{equation}
  xL_n'(x)=n\bigl[L_n(x)-L_{n-1}(x)\bigr],
  \qquad n\ge1.
  \label{eq:laguerre-deriv}
\end{equation}

To eliminate derivatives and obtain the three-term recurrence, first apply
\eqref{eq:laguerre-deriv} with index $n+1$:
\begin{equation}
  xL_{n+1}'=(n+1)(L_{n+1}-L_n).
  \label{eq:laguerre-deriv-nplus1}
\end{equation}
Next use \eqref{eq:laguerre-derivative-shift} with index $n+1$:
\[
  L_{n+1}'=L_n'-L_n.
\]
Multiplying by $x$ and then using \eqref{eq:laguerre-deriv} gives
\begin{align*}
  xL_{n+1}'
   &=xL_n'-xL_n\\
   &=n(L_n-L_{n-1})-xL_n.
\end{align*}
Equating this expression with \eqref{eq:laguerre-deriv-nplus1},
\[
  (n+1)(L_{n+1}-L_n)
    =nL_n-nL_{n-1}-xL_n.
\]
Expand the left-hand side and isolate $L_{n+1}$:
\begin{align*}
  (n+1)L_{n+1}-(n+1)L_n
    &=nL_n-nL_{n-1}-xL_n,\\
  (n+1)L_{n+1}
    &=(n+1)L_n+nL_n-xL_n-nL_{n-1},\\
  (n+1)L_{n+1}
    &=(2n+1-x)L_n-nL_{n-1}.
\end{align*}
Therefore
\begin{equation}
  (n+1)L_{n+1}(x)
    =(2n+1-x)L_n(x)-nL_{n-1}(x),
  \qquad n\ge1.
  \label{eq:laguerre-three-term}
\end{equation}
For $n=0$, the same formula reads
$L_1=(1-x)L_0$, which is also correct.  As a check at $n=1$,
\begin{align*}
  2L_2
   &=(3-x)L_1-L_0\\
   &=(3-x)(1-x)-1\\
   &=3-4x+x^2-1\\
   &=2-4x+x^2,
\end{align*}
so $L_2=1-2x+x^2/2$, in agreement with the direct series calculation.
\end{frame}

\begin{frame}{Useful formulas for the Laguerre polynomial}
\footnotesize
\begin{center}
\renewcommand{\arraystretch}{1.7}
\begin{tabular}{ll}
differential equation
  & $xL_n''+(1-x)L_n'+nL_n=0$ \\
series solution
  & $L_n(x)=\displaystyle\sum_{k=0}^{n}
       \frac{(-1)^k n!}{(n-k)!(k!)^2}x^k$ \\
generating function
  & $\dfrac{1}{1-t}\exp\!\left(-\dfrac{xt}{1-t}\right)
       =\displaystyle\sum_{n=0}^{\infty}L_n(x)t^n$ \\
derivative identities
  & $L_n'=L_{n-1}'-L_{n-1}$ \\
  & $xL_n'=n(L_n-L_{n-1})$ \\
three-term recurrence
  & $(n+1)L_{n+1}=(2n+1-x)L_n-nL_{n-1}$ \\
orthonormality
  & $\displaystyle\int_0^{\infty}L_m(x)L_n(x)e^{-x}\,dx=\delta_{mn}$ \\
Rodrigues
  & $L_n(x)=\dfrac{e^x}{n!}\dfrac{d^n}{dx^n}
       \left(x^ne^{-x}\right)$ \\
Schl\"afli
  & $L_n(x)=\dfrac{1}{2\pi i}
       \displaystyle\oint_{|z|=\rho}
       \frac{\exp\!\left(-xz/(1-z)\right)}
            {(1-z)z^{n+1}}\,dz,\quad 0<\rho<1$
\end{tabular}
\end{center}
\vspace{0.2em}
{\scriptsize The generating series requires $|t|<1$.
The derivative identities hold for $n\ge1$; the three-term recurrence extends
to $n=0$ using $L_0=1$.}
\end{frame}

% ======================================================================
\section{Chebyshev Polynomial}
% ======================================================================

\begin{frame}[allowframebreaks]{Definition through $\cos n\theta$}
\footnotesize

The Chebyshev polynomial of the first kind is defined by
\begin{equation}
  T_n(\cos\theta)=\cos(n\theta).
  \label{eq:cheby-def}
\end{equation}
That is, set $x=\cos\theta$ and read off a polynomial in $x$. The first cases:
\[
  T_0(x) = \cos 0 = 1,
  \qquad T_1(x) = \cos\theta = x .
\]
For $\cos 2\theta$,
\[
  \cos 2\theta = \cos^2\theta - \sin^2\theta = 2\cos^2\theta - 1 = 2x^2 - 1,
\]
so $T_2(x)=2x^2-1$.


For $\cos 3\theta$, use $\cos 3\theta+i\sin 3\theta
=e^{i3\theta}=(\cos\theta+i\sin\theta)^3$. Taking the real part,
\begin{align*}
  T_3(\cos\theta)
   &= \operatorname{Re}(\cos\theta+i\sin\theta)^3 \\
   &= \operatorname{Re}\bigl(\cos^3\theta + 3i\sin\theta\cos^2\theta
        - 3\sin^2\theta\cos\theta - i\sin^3\theta\bigr) \\
   &= \cos^3\theta - 3(1-\cos^2\theta)\cos\theta \\
   &= 4\cos^3\theta - 3\cos\theta .
\end{align*}
Therefore
\begin{equation}
  T_3(x) = 4x^3 - 3x .
  \label{eq:T3}
\end{equation}
\medskip
The name has many transliterations: Chebyshev, Tchebychev, Tschebyscheff, and more. The
mathematician's first name is less ambiguous: Pafnuty.
\end{frame}

\begin{frame}[allowframebreaks]{Recurrence and first few Chebyshev polynomials}
\footnotesize

The angle-addition identities are
\begin{align*}
  \cos\bigl((n+1)\theta\bigr)
    &=\cos(n\theta)\cos\theta-\sin(n\theta)\sin\theta,\\
  \cos\bigl((n-1)\theta\bigr)
    &=\cos(n\theta)\cos\theta+\sin(n\theta)\sin\theta.
\end{align*}
Adding them cancels the sine terms:
\begin{align*}
  \cos\bigl((n+1)\theta\bigr)
   +\cos\bigl((n-1)\theta\bigr)
    &=2\cos(n\theta)\cos\theta.
\end{align*}
Now set $x=\cos\theta$ and use
$T_j(\cos\theta)=\cos(j\theta)$:
\[
  T_{n+1}(x)+T_{n-1}(x)=2xT_n(x).
\]
Thus
\begin{equation}
  T_{n+1}(x)=2xT_n(x)-T_{n-1}(x),
  \qquad n\ge1,
  \label{eq:cheby-recur}
\end{equation}
with initial values $T_0(x)=1$ and $T_1(x)=x$.

For example, the recurrence gives
\begin{align*}
  T_2(x)
   &=2xT_1(x)-T_0(x)\\
   &=2x(x)-1\\
   &=2x^2-1,\\[0.4em]
  T_3(x)
   &=2xT_2(x)-T_1(x)\\
   &=2x(2x^2-1)-x\\
   &=4x^3-2x-x\\
   &=4x^3-3x,\\[0.4em]
  T_4(x)
   &=2xT_3(x)-T_2(x)\\
   &=2x(4x^3-3x)-(2x^2-1)\\
   &=8x^4-6x^2-2x^2+1\\
   &=8x^4-8x^2+1.
\end{align*}

For $n\ge1$, the leading coefficient of $T_n$ is $2^{n-1}$.  This follows by
induction.  The base case is $T_1=x$, whose leading coefficient is
$1=2^0$.  Suppose the leading term of $T_n$ is $2^{n-1}x^n$.  In
\eqref{eq:cheby-recur}, $T_{n-1}$ has degree $n-1$, so it cannot affect the
coefficient of $x^{n+1}$.  Therefore
\[
  T_{n+1}(x)
    =2x\left(2^{n-1}x^n+\text{lower powers}\right)
      -T_{n-1}(x)
    =2^n x^{n+1}+\text{lower powers}.
\]
This completes the induction.

The endpoint values follow directly from the trigonometric definition:
\begin{align*}
  T_n(1)
    &=T_n(\cos0)
      =\cos(n\cdot0)
      =1,\\
  T_n(-1)
    &=T_n(\cos\pi)
      =\cos(n\pi)
      =(-1)^n.
\end{align*}
The same definition also gives the parity relation
\begin{align*}
  T_n(-x)
    &=T_n\!\left(\cos(\pi-\theta)\right)\\
    &=\cos\!\left(n(\pi-\theta)\right)\\
    &=\cos(n\pi)\cos(n\theta)+\sin(n\pi)\sin(n\theta)\\
    &=(-1)^nT_n(x).
\end{align*}
\end{frame}

\begin{frame}[allowframebreaks]{Generating functions for $T_n$}
\footnotesize

Let
\[
  G(x,t):=\sum_{n=0}^{\infty}T_n(x)t^n.
\]
Multiply \eqref{eq:cheby-recur} by $t^{n+1}$ and sum over $n\ge1$:
\[
  \sum_{n=1}^{\infty}T_{n+1}t^{n+1}
    =2x\sum_{n=1}^{\infty}T_nt^{n+1}
     -\sum_{n=1}^{\infty}T_{n-1}t^{n+1}.
\]
Each sum can be expressed in terms of $G$:
\begin{align*}
  \sum_{n=1}^{\infty}T_{n+1}t^{n+1}
    &=\sum_{j=2}^{\infty}T_jt^j
      =G-T_0-T_1t
      =G-1-xt,\\
  2x\sum_{n=1}^{\infty}T_nt^{n+1}
    &=2xt\sum_{n=1}^{\infty}T_nt^n
      =2xt(G-T_0)
      =2xt(G-1),\\
  \sum_{n=1}^{\infty}T_{n-1}t^{n+1}
    &=t^2\sum_{j=0}^{\infty}T_jt^j
      =t^2G.
\end{align*}
Hence
\[
  G-1-xt=2xt(G-1)-t^2G.
\]
Collect the terms containing $G$:
\begin{align*}
  G-2xtG+t^2G
    &=1+xt-2xt\\
    &=1-xt.
\end{align*}
Therefore
\begin{equation}
  G(x,t)
    =\sum_{n=0}^{\infty}T_n(x)t^n
    =\frac{1-xt}{1-2xt+t^2}.
  \label{eq:cheby-gen-standard}
\end{equation}
For $x\in[-1,1]$, this analytic series converges for $|t|<1$; the identity
also holds as an identity of formal power series for arbitrary $x$.

A commonly used symmetric form weights all positive indices by $2$.  Since
\[
  T_0(x)+2\sum_{n=1}^{\infty}T_n(x)t^n
    =2G(x,t)-T_0(x)
    =2G(x,t)-1,
\]
we find
\begin{align*}
  2G-1
   &=
    \frac{2(1-xt)}{1-2xt+t^2}
    -\frac{1-2xt+t^2}{1-2xt+t^2}\\
   &=
    \frac{2-2xt-1+2xt-t^2}{1-2xt+t^2}\\
   &=
    \frac{1-t^2}{1-2xt+t^2}.
\end{align*}
Thus
\begin{equation}
  \frac{1-t^2}{1-2xt+t^2}
    =T_0(x)+2\sum_{n=1}^{\infty}T_n(x)t^n,
  \qquad x\in[-1,1],\quad |t|<1.
  \label{eq:cheby-gen-symmetric}
\end{equation}

As a coefficient check, expand the denominator equation
\[
  (1-2xt+t^2)G=1-xt.
\]
The coefficients of $t^0$ and $t^1$ give $T_0=1$ and $T_1=x$,
respectively.  For $t^{n+1}$ with $n\ge1$, the coefficient is
\[
  T_{n+1}-2xT_n+T_{n-1}=0,
\]
which recovers \eqref{eq:cheby-recur}.
\end{frame}

\begin{frame}{Shapes of the first Chebyshev polynomials}
\footnotesize
\begin{center}
  \includegraphics[width=0.86\linewidth]{figs/chebyshev-T.pdf}
\end{center}
The identity $T_n(\cos\theta)=\cos(n\theta)$ keeps every $T_n$ between $-1$ and $1$ on
the interval $[-1,1]$, with alternating extrema inherited from the cosine.
\end{frame}

\begin{frame}[allowframebreaks]{Chebyshev nodes}
\footnotesize
\begin{center}
  \includegraphics[width=0.86\linewidth]{figs/chebyshev-nodes.pdf}
\end{center}

The zeros of $T_n$ are obtained directly from
$T_n(\cos\theta)=\cos(n\theta)$.  On $x\in[-1,1]$, write
$x=\cos\theta$ with $0\le\theta\le\pi$.  Then
\begin{align*}
  T_n(x)=0
  &\iff \cos(n\theta)=0\\
  &\iff n\theta=\frac{(2k-1)\pi}{2},
       \qquad k=1,2,\ldots,n\\
  &\iff \theta_k=\frac{(2k-1)\pi}{2n}.
\end{align*}
Therefore the $n$ distinct zeros are
\begin{equation}
  x_k=\cos\!\left(\frac{(2k-1)\pi}{2n}\right),
  \qquad k=1,2,\ldots,n.
  \label{eq:cheby-nodes}
\end{equation}
They are distinct because
\[
  0<\theta_1<\theta_2<\cdots<\theta_n<\pi
\]
and $\cos\theta$ is strictly decreasing on $[0,\pi]$.

For example, when $n=4$,
\begin{align*}
  x_1&=\cos\frac{\pi}{8},&
  x_2&=\cos\frac{3\pi}{8},&
  x_3&=\cos\frac{5\pi}{8},&
  x_4&=\cos\frac{7\pi}{8}.
\end{align*}
Using $\cos(\pi-\alpha)=-\cos\alpha$, these occur in symmetric pairs:
\[
  x_4=-x_1,\qquad x_3=-x_2.
\]
The endpoint clustering can be quantified from the largest zero:
\begin{align*}
  1-x_1
   &=1-\cos\!\left(\frac{\pi}{2n}\right)\\
   &=2\sin^2\!\left(\frac{\pi}{4n}\right)
     \sim \frac{\pi^2}{8n^2}
     \qquad (n\to\infty).
\end{align*}
Thus the nearest node lies at distance $O(n^{-2})$ from each endpoint, rather
than at the $O(n^{-1})$ spacing characteristic of equally spaced nodes.  This
endpoint concentration is central to the stability of Chebyshev interpolation.
\end{frame}

\begin{frame}[allowframebreaks]{Differential equation and explicit series}
\footnotesize

Let
\[
  u(\theta):=T_n(\cos\theta)=\cos(n\theta).
\]
Since
\[
  \frac{d^2u}{d\theta^2}+n^2u=0,
\]
we convert the derivatives from $\theta$ to $x=\cos\theta$.  The chain rule gives
\[
  \frac{dx}{d\theta}=-\sin\theta,
  \qquad
  \frac{d}{d\theta}
    =\frac{dx}{d\theta}\frac{d}{dx}
    =-\sin\theta\,\frac{d}{dx}.
\]
Apply this operator a second time to a function $y(x)$:
\begin{align*}
  \frac{d^2y}{d\theta^2}
   &=
    \frac{d}{d\theta}
      \left(-\sin\theta\,\frac{dy}{dx}\right)\\
   &=
    -\cos\theta\,\frac{dy}{dx}
    -\sin\theta\,
      \frac{d}{d\theta}\left(\frac{dy}{dx}\right)\\
   &=
    -\cos\theta\,y'
    -\sin\theta
      \left(
        \frac{dx}{d\theta}\frac{d}{dx}y'
      \right)\\
   &=
    -\cos\theta\,y'
    -\sin\theta(-\sin\theta)y''\\
   &=
    -\cos\theta\,y'
    +\sin^2\theta\,y''.
\end{align*}
Using $\cos\theta=x$ and $\sin^2\theta=1-x^2$,
\[
  \frac{d^2y}{d\theta^2}
    =(1-x^2)y''-xy'.
\]
Therefore $y(x)=T_n(x)$ satisfies
\begin{equation}
  (1-x^2)T_n''(x)-xT_n'(x)+n^2T_n(x)=0.
  \label{eq:cheby-eq}
\end{equation}

We now derive the explicit finite series.  Let
\[
  y(x)=\sum_{j=0}^{\infty}a_jx^j.
\]
Then
\begin{align*}
  y'
    &=\sum_{j=0}^{\infty}(j+1)a_{j+1}x^j,\\
  y''
    &=\sum_{j=0}^{\infty}(j+2)(j+1)a_{j+2}x^j,\\
  -x^2y''
    &=-\sum_{j=2}^{\infty}j(j-1)a_jx^j
      =-\sum_{j=0}^{\infty}j(j-1)a_jx^j,\\
  -xy'
    &=-\sum_{j=1}^{\infty}j\,a_jx^j
      =-\sum_{j=0}^{\infty}j\,a_jx^j.
\end{align*}
Substitution into \eqref{eq:cheby-eq} gives
\begin{align*}
  0
   &=
    \sum_{j=0}^{\infty}
    \left[
      (j+2)(j+1)a_{j+2}
      -j(j-1)a_j-ja_j+n^2a_j
    \right]x^j\\
   &=
    \sum_{j=0}^{\infty}
    \left[
      (j+2)(j+1)a_{j+2}
      +(n^2-j^2)a_j
    \right]x^j.
\end{align*}
Hence
\begin{equation}
  a_{j+2}
    =-\frac{n^2-j^2}{(j+2)(j+1)}a_j.
  \label{eq:cheby-coeff-up}
\end{equation}
The even and odd coefficients are independent.  For the degree-$n$ polynomial,
only powers with the same parity as $n$ occur.  It is convenient to solve the
recurrence downward.  Rearranging \eqref{eq:cheby-coeff-up} and setting
$j=n-2k$ gives, for $k\ge1$,
\begin{align*}
  a_{n-2k}
   &=
    -\frac{(n-2k+2)(n-2k+1)}
            {n^2-(n-2k)^2}
      a_{n-2k+2}\\
   &=
    -\frac{(n-2k+2)(n-2k+1)}
            {n^2-\left(n^2-4nk+4k^2\right)}
      a_{n-2k+2}\\
   &=
    -\frac{(n-2k+2)(n-2k+1)}
            {4k(n-k)}
      a_{n-2k+2}.
\end{align*}
Iterating this relation $k$ times yields
\begin{align*}
  a_{n-2k}
   &=
    (-1)^k a_n
    \prod_{r=1}^{k}
      \frac{(n-2r+2)(n-2r+1)}
           {4r(n-r)}\\
   &=
    (-1)^k a_n
    \frac{n!/(n-2k)!}
         {4^k k!\,(n-1)!/(n-k-1)!}.
\end{align*}
For $n\ge1$, the leading coefficient derived from the three-term recurrence is
$a_n=2^{n-1}$.  Therefore
\begin{align*}
  a_{n-2k}
   &=
    (-1)^k 2^{n-1}
    \frac{n!}{(n-2k)!}
    \frac{(n-k-1)!}{4^k k!(n-1)!}\\
   &=
    (-1)^k
    \frac{n\,2^{n-2k-1}(n-k-1)!}
         {k!(n-2k)!}.
\end{align*}
Multiplying this coefficient by $x^{n-2k}$ and summing over all allowable $k$
gives
\begin{equation}
  T_n(x)
   =\frac{n}{2}
    \sum_{k=0}^{\lfloor n/2\rfloor}
      \frac{(-1)^k(n-k-1)!}
           {k!(n-2k)!}
      (2x)^{n-2k},
  \qquad n\ge1.
  \label{eq:cheby-explicit}
\end{equation}
The exceptional case is $T_0(x)=1$.

For $n=4$, formula \eqref{eq:cheby-explicit} gives
\begin{align*}
  T_4(x)
   &=
    2\left[
      \frac{3!}{0!\,4!}(2x)^4
      -\frac{2!}{1!\,2!}(2x)^2
      +\frac{1!}{2!\,0!}(2x)^0
    \right]\\
   &=
    2\left[
      \frac{6}{24}(16x^4)
      -\frac{2}{2}(4x^2)
      +\frac{1}{2}
    \right]\\
   &=
    2\left(4x^4-4x^2+\frac12\right)\\
   &=8x^4-8x^2+1,
\end{align*}
which agrees with the recurrence calculation.
\end{frame}

\begin{frame}[allowframebreaks]{Chebyshev orthogonality and norm}
\footnotesize

Multiply \eqref{eq:cheby-eq} by $(1-x^2)^{-1/2}$.  The first two terms become
\begin{align*}
  &(1-x^2)^{-1/2}
    \left[(1-x^2)T_n''-xT_n'\right]\\
  &\qquad=
    \sqrt{1-x^2}\,T_n''
    -\frac{x}{\sqrt{1-x^2}}T_n'.
\end{align*}
Since
\[
  \frac{d}{dx}\sqrt{1-x^2}
    =-\frac{x}{\sqrt{1-x^2}},
\]
the product rule gives
\[
  \frac{d}{dx}
    \left(\sqrt{1-x^2}\,T_n'\right)
   =
    \sqrt{1-x^2}\,T_n''
    -\frac{x}{\sqrt{1-x^2}}T_n'.
\]
Thus the differential equation has Sturm--Liouville form
\begin{equation}
  \frac{d}{dx}
    \left(\sqrt{1-x^2}\,T_n'\right)
  +\frac{n^2}{\sqrt{1-x^2}}T_n=0.
  \label{eq:cheby-sl}
\end{equation}

To verify orthogonality from this form, write the equations for $T_n$ and $T_m$:
\begin{align*}
  \left(\sqrt{1-x^2}\,T_n'\right)'
    +\frac{n^2}{\sqrt{1-x^2}}T_n&=0,\\
  \left(\sqrt{1-x^2}\,T_m'\right)'
    +\frac{m^2}{\sqrt{1-x^2}}T_m&=0.
\end{align*}
Multiply the first by $T_m$, the second by $T_n$, subtract, and integrate:
\begin{align*}
  &(n^2-m^2)
    \int_{-1}^{1}
      \frac{T_n(x)T_m(x)}{\sqrt{1-x^2}}\,dx\\
  &\qquad=
  -\left[
    \sqrt{1-x^2}
    \left(T_mT_n'-T_nT_m'\right)
  \right]_{-1}^{1}.
\end{align*}
The polynomials and their derivatives remain finite at $x=\pm1$, whereas
$\sqrt{1-x^2}\to0$.  Hence the boundary term is zero.  For nonnegative
indices $n\ne m$, one has $n^2\ne m^2$, so the integral is zero.

The norm and the same orthogonality relation can be evaluated explicitly.
Set $x=\cos\theta$.  As $x$ increases from $-1$ to $1$, $\theta$ decreases
from $\pi$ to $0$, and
\[
  dx=-\sin\theta\,d\theta,
  \qquad
  \sqrt{1-x^2}=\sin\theta
  \quad (0\le\theta\le\pi).
\]
Therefore
\begin{align*}
  \int_{-1}^{1}
    \frac{T_n(x)T_m(x)}{\sqrt{1-x^2}}\,dx
   &=
    \int_{\pi}^{0}
      \frac{\cos(n\theta)\cos(m\theta)}
           {\sin\theta}
      \left(-\sin\theta\,d\theta\right)\\
   &=
    \int_{0}^{\pi}
      \cos(n\theta)\cos(m\theta)\,d\theta.
\end{align*}

When $n\ne m$, use
$2\cos A\cos B=\cos(A-B)+\cos(A+B)$:
\begin{align*}
  \int_0^\pi\cos(n\theta)\cos(m\theta)\,d\theta
   &=
    \frac12
    \int_0^\pi
      \left[
        \cos((n-m)\theta)+\cos((n+m)\theta)
      \right]d\theta\\
   &=
    \frac12
    \left[
      \frac{\sin((n-m)\theta)}{n-m}
      +\frac{\sin((n+m)\theta)}{n+m}
    \right]_{\theta=0}^{\theta=\pi}\\
   &=0.
\end{align*}
If one index is zero, the same conclusion follows directly from
$\int_0^\pi\cos(k\theta)\,d\theta=0$ for every positive integer $k$.

When $n=m\ge1$,
\begin{align*}
  \int_0^\pi\cos^2(n\theta)\,d\theta
   &=
    \frac12\int_0^\pi
      \left[1+\cos(2n\theta)\right]d\theta\\
   &=
    \frac12
    \left[
      \theta+\frac{\sin(2n\theta)}{2n}
    \right]_{0}^{\pi}\\
   &=\frac{\pi}{2}.
\end{align*}
For $n=m=0$,
\[
  \int_0^\pi\cos^2(0\theta)\,d\theta
   =\int_0^\pi1\,d\theta
   =\pi.
\]
Combining all cases,
\begin{equation}
  \int_{-1}^{1}
    \frac{T_n(x)T_m(x)}{\sqrt{1-x^2}}\,dx
  =
  \begin{cases}
    \pi, & n=m=0,\\[0.2em]
    \pi/2, & n=m\ge1,\\[0.2em]
    0, & n\ne m.
  \end{cases}
  \label{eq:cheby-ortho}
\end{equation}
\end{frame}

\begin{frame}{Useful formulas for the Chebyshev polynomial}
\footnotesize
\begin{center}
\renewcommand{\arraystretch}{1.65}
\begin{tabular}{ll}
differential equation
  & $(1-x^2)T_n''-xT_n'+n^2T_n=0$ \\
explicit series ($n\ge1$)
  & $T_n(x)=\dfrac{n}{2}
       \displaystyle\sum_{k=0}^{\lfloor n/2\rfloor}
       \dfrac{(-1)^k(n-k-1)!}{k!(n-2k)!}(2x)^{n-2k}$ \\
standard generating function
  & $\dfrac{1-xt}{1-2xt+t^2}
       =\displaystyle\sum_{n=0}^{\infty}T_n(x)t^n$ \\
symmetric generating function
  & $\dfrac{1-t^2}{1-2xt+t^2}
       =T_0(x)+2\displaystyle\sum_{n=1}^{\infty}T_n(x)t^n$ \\
recurrence
  & $T_{n+1}=2xT_n-T_{n-1}$ \\
zeros
  & $x_k=\cos\!\left(\dfrac{(2k-1)\pi}{2n}\right),
       \quad k=1,\ldots,n$ \\
orthogonality and norm
  & $\displaystyle\int_{-1}^{1}
       \dfrac{T_nT_m}{\sqrt{1-x^2}}\,dx
       =\begin{cases}
          \pi, & n=m=0,\\
          \pi/2, & n=m\ge1,\\
          0, & n\ne m .
        \end{cases}$
\end{tabular}
\end{center}
\vspace{0.2em}
{\footnotesize For $x\in[-1,1]$, both generating series converge for
$|t|<1$.}
\end{frame}

\begin{frame}{Summary: four orthogonal families}
\footnotesize
\begin{center}
\renewcommand{\arraystretch}{1.5}
\begin{tabular}{lllc}
\textbf{Family} & \textbf{Interval} & \textbf{Weight $w(x)$} & \textbf{Arises in} \\
\hline
Legendre $P_n$ & $[-1,1]$ & $1$ & spherical harmonics \\
Hermite $H_n$  & $(-\infty,\infty)$ & $e^{-x^2}$ & quantum oscillator \\
Laguerre $L_n$ & $[0,\infty)$ & $e^{-x}$ & hydrogen ($L_n^{(\alpha)}$) \\
Chebyshev $T_n$ & $[-1,1]$ & $(1-x^2)^{-1/2}$ & $\cos n\theta$, interpolation \\
\end{tabular}
\end{center}
\vspace{0.6em}
The shared structure is a second-order Sturm--Liouville ODE, a generating function, a
recurrence, and an orthogonality relation on the natural interval and weight. For Legendre,
Hermite, and Laguerre we also recorded Rodrigues' formulas and Schl\"afli contour integrals.
Completeness then lets us expand reasonable functions in the corresponding series, exactly
as Fourier series expand in $e^{inx}$.
\end{frame}

\end{document}
